% ============================================================
%  第九卷 第七章 磁性（§69–§74*）
%  底本：高教社中文第4版（王锡绂 译）；印刷页 265–296（PDF p276–p307）
%  带 \chapter（章名"磁性"，章名页 p276/印刷页265）；\setcounter{section}{68}
%  保证单独编译时节号/公式号与原书一致（全书连排时 ch06b 已至 §68，幂等无害）。
%
%  插图清单：本章无插图（原书图 10–18 均不在本章范围；§70 自旋波谱无小图，
%  p282–p284 经核对无任何插图）。
%
%  圈码脚注：正文圈码共 28 处原位保留（p276①、p277①、p278①、p279①、
%  p281①、p282①②、p284①②、p285①、p286①②③、p288①、p290①、
%  p291①②、p292①、p294①、p295①、p297①、p298①、p299①、p302①②、
%  p303①、p304①、p305①）。脚注文本已全部按原书页面 300dpi 页脚裁片逐字转录
%  为 \footnote{...}（脚注内黑体强调用 \textbf、西里尔人名照录）。
%  其中 p278①（"自由能 F̃①"）、p282②（"能谱 ε=ℏω②"）、p284①
%  （"能隙的出现①"）、p295①（"χ_mn=χ_nm①"）、p302①（"正比于 k²①"）
%  五处 OCR 丢失或残缺，均按页面补；p275 为第六章末页，其①②与本任务无关。
%
%  编号公式核算：OCR 带 tag 103 处（其中 74*.4、74*.6 误作 74°.4、74°.6，
%  74*.8 误作 74.8、74*.14 误作 74.14，均照页面归位），与页面逐条核对：
%  §69 (69.1–69.14) 14、§70 (70.1–70.12) 12、§71 (71.1–71.13) 13、
%  §72 (72.1–72.17) 17、§73 (73.1–73.14) 14、§74 (74.1–74.19) 19、
%  §74* (74*.1–74*.14) 14，共 103 个编号公式，无重号、无缺号、无跳号。
%  另有未编号展示公式约 60 个照页面排 equation*。
%
%  已修正的 OCR 错误（均对照 原书转换/pages/pNNN.png 逐条核实）：
%   1. 页锚错位：OCR 首锚 %%p 275 实为第七章章名页 p276（p275 为第六章
%      末页）；本文件 %%p 锚已按实际页面重新标注。
%   2. 汉字上下标按页面还原：H_有效、U_非匀（p279，OCR U_{\ast\ast\jmath}/
%      U_{\F'\Game}）、U_异（p280–p304，OCR U_{☉}/U_{\Re}/U_{\mathfrak H}）、
%      Ĥ_交换（p276/p290，OCR Ĥ_{☉*}/Ĥ_{\pm\rlap/\Delta x}）、F_交换
%      （p301，OCR F_{\dot x\oplus\vec x}）、ε_交换（p288）、C_磁
%      （p295，OCR C_{ss}）、Ω_互、M_互、C_互（p298–p299，OCR Ω_Ξ/M_π/C_x）。
%   3. p281：(69.14) 前变分式按页面为 −KM_zν·δM（OCR ν·{δM 乱码）。
%   4. p282：(70.5) 下标按页面为 x,y,z（OCR M_s/M_γ/M_s）；(70.6) 右端
%      δ(r₁−r₂)（OCR S(r₁−r₂)）；"能谱 ε=ℏω②"补脚注号②。
%   5. p283："（这里写的方程已乘过因子 ℏ）"（OCR 因子 h）；(70.10) 后
%      "sin²θ=k_x²/K²"照页面（见排印问题 c）；δU_异=K′M²ϑ²（OCR \mathfrak H）。
%   6. p284 习题：方程按页面为 −iεm=2β{αk²n_x×m−|K|m_xn_y−n_x×h}；
%      "所有 k① 的方向"圈码归位（OCR k\textcircled{1} 游离）。
%   7. p286："（下标用 mag 表示——译注）"照页面补入；(71.5) 常数照页面
%      0.113（见排印问题 p）；(71.7) 按页面排两行 aligned。
%   8. p288：(71.11) ε_交换 下标（OCR ε_{\tilde x\tilde x}）；∂δM/∂M
%      被积式照页面；(71.13) 分母 15α^{3/2}𝔥^{1/2}（OCR ζ_Q^{1/2}）。
%   9. p290 习题4：⟨â_k⁺â_k⟩=n_k（OCR nk）；φ_ik 各式照页面。
%  10. p291：(72.1) 按页面为 J_nm Ŝ_nŜ_m（OCR Ŝ_nŜ_m 顺序同）；"格点的
%      径矢 r_n"（OCR r_π）。
%  11. p292：(72.9) (2S)^{−1/2}Ŝ_n₋χ₀；p292 脚注锚"减少 1①"（OCR 使-
%      一个粘连）；"NS−1 值"照页面。
%  12. p293：(72.11) 指数 e^{ik·r_n}；(72.12) 求和号带撇 Σ_m′（OCR et'
%      乱码）；(72.14) χ_qiχ_qk 照页面（原书沿用俄文 х 记 r_q 之分量，
%      中译本排作 χ，照排）。
%  13. p294：Z 展开式末项 1/6T³Ē³（放大核对为 T³）；"现在 Z 的对数再按
%      1/T 的幂展开"（OCR \stackrel E=0 乱码窜入）；F 式 \overline{S²_{nx}}
%      照页面（见排印问题 r）。
%  14. p295：(73.1) χ_nn 归一化因子 [4S(2S−1)]^{−1/2}；"χ_mn=χ_nm①"
%      圈码保留。
%  15. p296：(73.5) 第二行按页面为 ψ_mm+ψ_nn（OCR ψ_mn）；(73.4) 末项
%      4β𝔥Ŝ_m₋Ŝ_n₋χ₀（OCR 4βŜŜ 乱码）；"对角"的 J_ll=0（OCR J_u）。
%  16. p297：(73.6) 指数 e^{iK·(r_m+r_n)/2}；p297 首行"再引进记号 J"。
%  17. p298：(73.10) 两行并一式（原书跨页断行）；(73.13) Ω_互=VT⁵f(𝔥/T)
%      （OCR T^s）。
%  18. p299：习题 M_互/M、C_互 照页面；ε(k)=SJ₀a²k²+2β𝔥。
%  19. p300：(74.3) 照页面 ∫Ωδℏ·SdV（见排印问题 n）；(74.5) K=L×H_L−
%      (1/γ)M×Ω。
%  20. p301：(74.9) H_L=(1/2)α_ik∂²L/∂x_i∂x_k, Ω=γaM；磁矩流张量
%      Π_il=γα_lk(L×∂L/∂x_k)_i（下标 il/lk 照页面）。
%  21. p302：(74.10) 第二方程照页面 ∂M/∂t=γL₀ν×H_L（见排印问题 s）；
%      "上面已用过相对小的 m."断句（OCR m·粘连）；"垂于 ν 的平面"照页面
%      （见排印问题 t）。
%  22. p303：(74.15) ε(0)=ℏγL₀√(aK)；(74.17) 分母 (α₂α₁²)^{3/2}ℏ³；
%      (74.19) 照页面 (2p)³（见排印问题 g）；L(T)−L(0) 式照页面。
%  23. p304：习题解"相垂的平面"照页面（见排印问题 u）；§74* 首段按页面
%      理顺"当交换积分 J_nm 为负值时，显然要出现反铁磁性""此表式对经典
%      矢量 S_n 简单取极小"（OCR J/s_\ast 游离错位）；(74*.1) 前按页面补
%      "表示成遵守玻色对易规则的算符. 将自旋 S 写成如下的自旋算符"。
%  24. p305：(74*.1) Ŝ₋ 带 â⁺、Ŝ₊ 带 â 照页面；{Ŝ_z,Ŝ₊} 例式照页面；
%      "其投影 S_z=−1①"圈码归位；"（â⁺â=2S）"（OCR a^\dagger a 混排）。
%  25. p306：(74*.4) b 项第二算符照页面 â（见排印问题 k）；(74*.5) 三行
%      一编号，J_k^{(1)} 定义式求和上限照页面 k（见排印问题 m）；(74*.6)
%      上标均为 † 号（OCR 误作 *）。
%  26. p307："准动量为 ℏk 的自旋波量子"照页面补 ℏ（OCR 丢）；⟨S_z⟩≈S−0.08
%      （OCR S\dot−）；末段断句照页面（见排印问题 q）。
%  27. 全章：\pmb→\bm；𝔥 一律 \mathfrak{H}；热力学势/角速度 Ω 用 \varOmega
%      （与 ch03b/ch04 一致）；δ 变分记号照排；句读统一半角"."；删除零散
%      OCR 噪声（"22 D""2 D""X D""8 D"等页码水印残迹、p292 末"22 D"）。
%
%  原书排印问题（照排并登记，供总校裁定）：
%   a. p281："因顾及相对论效应而导致运动方程(69.5)发生的变化"，按上下文
%      疑应为 (69.4)，照排。
%   b. p283："麦克斯韦方程(69.10)给出"，按上下文（准静场方程）疑应为
%      (69.9)，照排。
%   c. p283："sin²θ=k_x²/K²"，分母 K 按物理意义应为波数 k（sin²θ=k_x²/k²），
%      照排。
%   d. p286："由此得到 αk≪1"，疑应为 αk²≪1，照排。
%   f. p299 习题："(72.12)式的积分计算导致以下的结果"，疑应为 (73.12)，
%      照排。
%   g. p303：(74.19) 分母照页面印作 (2p)³，疑应为 (2π)³，照排。
%   h. p303："在(74.12)式中将 αk² 换成 αk²+K 即可得到自旋波的色散律"，
%      疑应为 (74.13)，照排。
%   i. p303："可以使用能谱的(74.12)式"，疑应为 (74.13)，照排。
%   j. p304 习题：方程 −iωm=−γLα(n)k²(l×ν)−γL|K|(n_z·l)n_y 照页面
%      （两项均负号）；按推导（H_L=−α(n)k²l−|K|(n_z·l)n_z 代入 ∂m/∂t=
%      γL₀ν×H_L）及该解的色散结论，两项均应为正号，照排。
%   k. p306：(74*.4) 照页面印作 b̂_a⁺â_a、b̂_a⁺â_a′，按 (74*.2)(74*.3)
%      及后文对角化疑应为 b̂_a⁺b̂_a、b̂_a⁺b̂_a′，照排。
%   l. p306："为此如(72.7)，设"，疑应为 (73.7)，照排。
%   m. p306：(74*.5) J_k^{(1)} 定义式求和号上限照页面印作 k，按上下文疑应
%      为 a′（对亚晶格节点求和），照排。
%   n. p300：(74.3) 左中项照页面印作 ∫Ωδℏ·SdV，按上下文疑应为
%      ∫Ω·δ(ℏS)dV，照排。
%   o. p307："准动量为 ℏk 的自旋波量子"，与 §70"准动量 k"体例不一，
%      照页面照排。
%   p. p286：(71.5) 数值常数照页面 0.113；按 5Γ(5/2)ζ(5/2)/8π²≈0.110
%      疑为排印（或计算）误差，照排。
%   q. p307 末段："在这类晶格对远邻有重要的. 与最近邻有不同符号的相互
%      作用"，断句/连接词疑有脱漏（"重要的"后疑应为顿号），照页面照排。
%   r. p294：F 式第一项 \overline{S²_{nx}} 照页面（下标 x 与随后平均值式
%      \overline{S²_{nz}}=S(S+1)/3 并用，数值相同不影响结果），照排。
%   s. p302：(74.10) 第二方程照页面 ∂M/∂t=γL₀ν×H_L，按上下文（线性化）
%      疑应为 ∂m/∂t，照排。
%   t. p302："但 l 在垂于 ν 的平面上"，疑脱"直"字，照排。
%   u. p304 习题："处于与晶体对称轴(z 轴)相垂的平面"，疑脱"直"字，照排。
%   v. p281 (69.14) 前变分项照页面 −KM_zν·δM（K 与 M_z 并写）。
%
%  遗留问题：
%   - 28 处圈码脚注文本已全部转录完毕（脚注代理按 pNNN 页脚补录，与页面逐字核对）；
%   - OCR 首页锚（%%p 275）实为 p276，脚注代理引用页码时以本文件锚为准。
% ============================================================

\setcounter{section}{68}

\chapter{磁性}

%%p 276
\section{铁磁体中的磁矩运动方程}

晶体的磁性结构使其能谱出现特殊的分支，在研究这些能谱时，我们首先回忆磁性物质中相互作用的一些特性.

铁磁体内相互作用的基本形式是原子的交换作用，由其导致自发磁化. 这种相互作用的特点在于它与磁化对晶格的取向无关：在考虑到系统波函数的对称性，交换相互作用是电子静电相互作用的结果，并与总自旋\footnote{对于铁磁体，旋磁比 $g$ 的实验值极接近数值 2，这一点证实了铁磁性的自旋本质.}的方向无关.

最简单的铁磁系统是晶格中的原子具有磁矩的电介质，并且磁矩平行时恰好在能量上是有利的，以此来确定交换相互作用的符号（“铁磁化”），这时，系统基态是所有的自旋平行的状态．确切地说，处于基态时系统的总自旋在某一方向的投影等于最大的可能值 $\sum s_a$（对所有原子求和），其中 $s_a$ 是一个原子的自旋. 事实上，交换作用的哈密顿量 $\hat{H}_{\text{交换}}$ 与系统的总自旋算符 $\hat{\bm{S}}$ 可对易，也就是与投影算符 $\hat{S}_z$ 可对易（因为 $\hat{H}_{\text{交换}}$ 与各个自旋的方向无关，而算符 $\hat{\bm{S}}$ 是自旋空间中的转动算符）. 因此基态应有确定的 $S_z$ 值，而能量的最小值对应于 $S_z$ 的最大值. 我们指出，这时每个原子自旋的投影 $s_z$ 等于其最大值 $s_a$，因此基态的磁矩等于自身的“额面”值 $\sum\mu_a$，其中 $\mu_a$ 是一个原子的磁矩. 但是，这个性质被更弱的、相对论性相互作用所破坏.

复杂一些的情况是物体的磁化强度不等于额面值. 尤其，不是所有原子间的相互作用都是铁磁性的时候，可能由两个相反磁化的亚晶格形成结构，亚晶格磁化强度不同，因此不能完全抵消；具有这种结构的物质称为铁氧体（完全抵消的情况，称为反铁磁体）.

%%p 277
最后，在研究铁磁性金属的原子自旋时，必然涉及传导电子，传导电子无论如何，甚至在 $T=0$ 时，也不可能完全磁化（由于费米简并效应）. 磁性相互作用的特性还使铁磁体具有更复杂的基态结构：磁矩原子的非共线分布——所谓螺旋结构.

与所有宏观系统一样，铁磁体的弱激发态可以看成元激发的集合，即准粒子气体. 原子磁矩有序分布中的元激发称为自旋波量子．既然所讨论的是平移对称晶格中的准粒子，于是，自旋波量子具有的动量就不是真实的动量，而是准动量，准动量取遍一个倒格胞中的所有值．在经典图像中，自旋波量子对应于自旋波，就是磁矩振动沿晶格的传播．自旋波量子遵从玻色统计，因此，自旋波量子态大的占有数对应于自旋波的经典极限情况.

如自旋波的波长大于晶格常数 $a$（即波矢 $k\ll1/a$），那么，这个波应看成是宏观的；于是，波的色散律 $\bm{\omega}(\bm{k})$ 将通过磁矩宏观运动方程中的唯象参数（物质常数）来表达. 因而自旋波量子的谱 $\bm{\varepsilon}=\hbar\bm{\omega}(\bm{k})$ 也将通过这些参数来表达. 有一种确定自旋波量子谱的方法，完全类似于用声波振动宏观方程中的宏观参数（弹性模量）来确定长波声子谱．为实现这一设想，必需预先导出上述运动方程\footnote{本节进一步的结果属于朗道和栗弗席兹的工作（1935）. 我们指出，这些结果对于“交换的”铁磁体是正确的. 在这里我们不涉及所谓的弱铁磁体，在这样的铁磁体中只在考虑相对论性效应时才显现出铁磁矩.}.

从只考虑交换相互作用的情况，来开始我们的研究.

我们感兴趣的是铁磁体的弱激发态（只有弱激发态的性质才可用普遍的形式加以解释），所以应该限于磁矩的低频“慢”运动．磁矩方向在空间缓慢改变、而磁矩大小保持不变的就是这样的运动. 事实上，磁化强度的平衡值是由交换相互作用确定的；因此，磁化强度的变化在任何波长都不与有限能耗相联系（假设，物体处于远离居里点的状态，在居里点没有自发磁化强度）. 另一方面，当只计及交换相互作用，物体作为整体，其磁矩转动时能量不变；因此波长越大，磁化强度在物体各处非均匀转动所需的能量就越小. 换句话说，长波振动的频率小.

决定磁矩方向变化的运动方程乃是在力矩 $\bm{K}$ 作用下自旋角动量 $\hbar\bm{S}$ 的进动方程：
\begin{equation*}
\frac{\partial\hbar\bm{S}}{\partial t}=\bm{K},
\end{equation*}
（见第一卷(34.4)）$\hbar\bm{S}$ 和 $\bm{K}$ 分别是单位体积的角动量和力矩. 磁矩密度（磁化强度）$\bm{M}$ 通过 $\bm{S}$ 依下式表达
\begin{equation*}
\bm{M}=\gamma\hbar\bm{S}=-\frac{g\,|e|}{2mc}\,\bm{S},
\end{equation*}

%%p 278
式中 $\gamma$ 是磁矩与角动量之比. $e=-|e|$ 和 $m$ 分别是电子的电荷和质量. $g$ 为铁磁体的旋磁比（对比第二卷§45）. 因此，$\bm{M}$ 的方程如下
\begin{equation}
\frac{\partial\bm{M}}{\partial t}=\gamma\bm{K}.
\end{equation}
如力矩 $\bm{K}$ 不依赖坐标，此方程为精确的. 如 $\bm{K}$ 在空间缓慢变化，此方程为近似的.

力矩 $\bm{K}$ 定义为自旋系统在无穷小转动时能量的变化. 令 $\delta\bm{\phi}$ 为沿旋转轴方向的矢量，其值等于转角，根据力学公式，力矩由 $\dfrac{\partial E}{\partial\bm{\phi}}$ 给出（见第一卷(34.6)）. 如 $\delta E$ 与坐标有关，应该将空间每点的微分换成变分．最好这个变分是在温度、体积以及物体内每点磁场都恒定的情况下进行，对于这些变量的热力学势是自由能 $\tilde{F}$\footnote{参阅第八卷§36（物体整体的热力学量，那里用草体拉丁字表示）. 在非均匀分布时，正确的说法是物体的自由能（当体积给定时），而不是热力学势 $\delta\tilde{\varPhi}$. 在这里，我们不关心磁致伸缩效应，即不考虑磁化强度变化时的应力和晶体的形变. 此时可不区分 $\delta\tilde{\varPhi}$, $\delta\tilde{F}$.}. 于是，力矩密度由以下方程确定：
\begin{equation}
\delta\tilde{F}=-\int\bm{K}\cdot\delta\bm{\phi}\,\mathrm{d}V.
\end{equation}
写出磁矩变化时（准确地说是方向变化，因其值固定）自由能的变化 $\delta\tilde{F}$ 为
\begin{equation}
\delta\tilde{F}=-\int\bm{H}_{\text{有效}}\cdot\delta\bm{M}\,\mathrm{d}V,
\end{equation}
此处类比外磁场内磁矩能的表达式而引入“有效”场 $\bm{H}_{\text{有效}}$. 平衡时 $\bm{H}_{\text{有效}}=0$. 因为磁矩的平衡分布刚好决定于自由能的极小条件. 现在考虑无穷小转动 $\delta\bm{\phi}$ 时，$\bm{M}$ 的变化等于
\begin{equation*}
\delta\bm{M}=\delta\bm{\phi}\times\bm{M}.
\end{equation*}
将此式代入(69.3)式并对比(69.2)式，得出力矩密度
\begin{equation*}
\bm{K}=\bm{M}\times\bm{H}_{\text{有效}}.
\end{equation*}
代入(69.1)式，最后得出磁矩的朗道–栗弗席兹方程：
\begin{equation}
\frac{\partial\bm{M}}{\partial t}=\frac{g\,|e|}{2mc}\,\bm{H}_{\text{有效}}\times\bm{M}.
\end{equation}
现在来证明，此方程表述的磁矩变化并不伴随能耗. 此能耗等于：
\begin{equation*}
Q=T\frac{\partial S}{\partial t}=-\frac{\partial R_{\min}}{\partial t}=-\frac{\partial\tilde{F}}{\partial t},
\end{equation*}
式中 $S$ 为物体的熵. $R_{\min}$ 为最小功，它是引导物体进入既定的非平衡态所必需的. 因此，利用(69.4)式，有

%%p 279
\begin{equation*}
Q=\int\bm{H}_{\text{有效}}\cdot\frac{\partial\bm{M}}{\partial t}\,\mathrm{d}V
=\frac{g\,|e|}{2mc}\int\bm{H}_{\text{有效}}\cdot
(\bm{H}_{\text{有效}}\times\bm{M})\,\mathrm{d}V=0,
\end{equation*}
于是论断得证. 应该指出，无能耗是靠方程(69.4)的右侧垂于 $\bm{H}_{\text{有效}}$ 来保证的.（我们在下节末还要讨论能耗问题）.

根据定义式(69.2)，对物体总自由能进行变分，可求出有效场的显式. 总自由能由如下积分给出
\begin{equation}
\tilde{F}=\int\Bigl[f_0(\bm{M})+U_{\text{非匀}}-\bm{M}\cdot\bm{H}-\frac{H^2}{8\pi}\Bigr]\mathrm{d}V
\end{equation}
（参阅第八卷§39）. 这里 $f_0(\bm{M})$ 是 $H=0$ 时均匀磁化体的自由能密度，因为只考虑交换相互作用，所以与 $\bm{M}$ 的方向无关. $U_{\text{非匀}}$ 是与矢量 $\bm{M}$ 的沿非均匀磁化体缓慢变化方向有关的附加交换能密度.

将此能量按磁矩 $\bm{M}$ 对坐标的微商的幂进行展开，头几项的形式为
\begin{equation}
U_{\text{非匀}}=\frac{1}{2}\alpha_{ik}
\frac{\partial\bm{M}}{\partial x_i}\frac{\partial\bm{M}}{\partial x_k}.
\end{equation}
并且，这个微商二次型本质上是正定的. (69.6)式的组成（根据交换作用的性质）应与矢量 $\bm{M}$ 的绝对方向无关. 单轴晶体的二阶对称张量 $\alpha_{ik}$ 具有分量 $\alpha_{xx}=\alpha_{yy}\equiv\alpha_1$，$\alpha_{zz}\equiv\alpha_2$（$z$ 轴是晶体的对称轴）；在立方晶体中 $\alpha_{ik}=\alpha\delta_{ik}$.

如果在数量级为晶格常数 $a$ 的距离上，磁矩方向有显著变化，再注意到晶格一个原胞非均匀性的能量应该达到交换相互作用能的原子特征值，就可以估计出系数 $\alpha_{ik}$ 的数量级．特征的交换能同居里温度 $T_{\mathrm{c}}$（铁磁性消失的温度）同数量级. 由条件 $T_{\mathrm{c}}/a^3\sim\alpha M^2/a^2$，我们得到
\begin{equation}
\alpha\sim T_{\mathrm{c}}/aM^2.
\end{equation}
在物体每一点都给定 $\bm{H}$ 时，对积分(69.5)进行变分，并对第二项进行分部积分，我们得到
\begin{equation*}
\delta\tilde{F}=\int\Bigl[f_0'(\bm{M})\frac{\bm{M}}{M}
-\alpha_{ik}\frac{\partial^2\bm{M}}{\partial x_i\partial x_k}
-\bm{H}\Bigr]\delta\bm{M}\,\mathrm{d}V.
\end{equation*}
根据定义(69.2)，花括号内的表达式是 $-\bm{H}_{\text{有效}}$. 第一项沿 $\bm{M}$ 的方向，但是，当代入运动方程(69.4)中时，第一项总会等于零，因此完全可以把它舍去\footnote{但是，此后在平衡时 $\bm{H}_{\text{有效}}$ 已不必变为零.}. 于是，我们得到
\begin{equation}
\bm{H}_{\text{有效}}=\alpha_{ik}\frac{\partial^2\bm{M}}{\partial x_i\partial x_k}+\bm{H}.
\end{equation}
为了得到完整的方程组，给(69.4)，(69.8)两式还要附加一个将场 $\bm{H}$ 与磁化强度 $\bm{M}$ 的分布相联系的麦克斯韦方程. 在下一节将要研究的自旋波是在 $\omega\ll ck$ 的意义上的低频波，在这些条件下的场是准静场；麦克斯韦方程中可以忽略对时间的微商，于是方程组变成，
\begin{equation}
\nabla\times\bm{H}=0,\qquad\nabla\cdot\bm{B}=\nabla\cdot(\bm{H}+4\pi\bm{M})=0.
\end{equation}

%%p 280
与此相关，在 $\bm{H}$ 恒定时，可能产生关于积分式(69.5)对 $\bm{M}$ 变分是否合理的问题，虽然 $\bm{M}$ 和 $\bm{H}$ 通过(69.9)的第二个方程相联系. 但是，如果取 $\bm{H}=-\nabla\varphi$（由于第一个方程）并计算积分对 $\varphi$ 的变分，那么由于第二个方程此变分为零，所以 $\bm{H}$ 的变分对 $\delta\tilde{F}$ 无贡献.

如果物体不在外磁场中，那么它内部的场便完全与磁化强度的分布有关，一般说来这个场是与 $\bm{M}$ 同一数量级的量. 在此意义上，有效场(69.8)中的 $\bm{H}$ 项是相对论效应（我们指出，原子磁矩以及和它有关的自发磁化强度都用玻尔磁子 $\beta=|e|\hbar/2mc$ 定义，这是一个在分母中含 $c$ 的量）. 因此，目前所研究的纯交换近似中，(69.8)的第二项应该略去，因此运动方程为
\begin{equation}
\frac{\partial\bm{M}}{\partial t}=\frac{g\,|e|}{2mc}\,\alpha_{ik}
\frac{\partial^2\bm{M}}{\partial x_i\partial x_k}\times\bm{M}.
\end{equation}
我们注意到这个方程是非线性的.

方程(69.10)可以写成磁矩连续性方程的形式：
\begin{equation}
\frac{\partial M_i}{\partial t}+\frac{\partial\varPi_{il}}{\partial x_l}=0,
\end{equation}
其中磁矩流密度张量为
\begin{equation*}
\varPi_{il}=\frac{g\,|e|}{2mc}\,\alpha_{ik}
\Bigl(\bm{M}\times\frac{\partial\bm{M}}{\partial x_k}\Bigr)_{\!i}.
\end{equation*}
这是预先应该想到的，因为在交换近似下物体的总磁矩守恒.

现在我们考虑，在铁磁体中除交换作用之外，在电子磁矩之间还有更弱的相对论性相互作用：自旋–自旋作用和自旋–轨道作用．在宏观理论中这些作用用各向异性磁能来描述，能量密度 $U_{\text{异}}$ 依赖于磁化强度矢量对晶格的方向；用这些相互作用来建立铁磁体的自发磁化强度的平衡方向. 上边已经指出过，$\bm{M}$ 同磁场 $\bm{H}$ 的相互作用也属于相对论性的.

在单轴晶体中各向异性的能量为
\begin{equation}
U_{\text{异}}=-\frac{K}{2}M_z^2.
\end{equation}
如果 $K>0$，那么平衡磁化方向就沿着对称轴——$z$ 轴（“易磁化轴”型铁磁体）的方向；如果 $K<0$，那么自发磁化方向就位于 $xy$ 平面上（“易磁化面”型铁磁体）. 立方晶体的各向异性能量可表为
\begin{equation}
U_{\text{异}}=\frac{K'}{M^2}
(M_x^2M_y^2+M_x^2M_z^2+M_y^2M_z^2).
\end{equation}
其中 $x,y,z$ 轴的方向沿着三个四级对称轴（立方晶胞的各棱边）．如果 $K'>0$，平衡矢量 $\bm{M}$ 的方向沿着立方晶胞的一个棱，如果 $K'<0$——则沿着晶胞空间对角

%%p 281
线中的一条\footnote{对于不同铁磁体量纲为 1 的量 $K,K'$ 在几十分之一到几十这个宽度范围内取值. 相对论相互作用与交换相互作用之比与 $a^3U_{\text{异}}/T_{\mathrm{c}}$ 同数量级，通常为 $10^{-4}\sim10^{-5}$.}.

为确定起见，我们将研究单轴铁磁体. 在(69.5)的被积分式中补充一项 $U_{\text{异}}$ (69.12)，对附加项变分后得到一项 $-KM_z\,\bm{\nu}\cdot\delta\bm{M}$，其中 $\bm{\nu}$ 为晶体对称轴方向上的单位矢量. 于是，对于有效场我们得到
\begin{equation}
\bm{H}_{\text{有效}}=\alpha_{ik}\frac{\partial^2\bm{M}}{\partial x_i\partial x_k}
+KM_z\bm{\nu}+\bm{H}.
\end{equation}
容易看出，有效场的这个改变已经概括了因顾及相对论效应而导致运动方程(69.5)发生的变化. 事实上，不存在耗散，运动方程的右侧仍然应当垂直于 $\bm{H}_{\text{有效}}$，即应当有 $\bm{M}'\times\bm{H}_{\text{有效}}$ 形式，其中 $\bm{M}'$ 和 $\bm{M}$ 的差别只能是相对论修正，这个修正值总远小于 $\bm{M}$，因而并不重要. $\bm{H}_{\text{有效}}$ 中的相对论各项附加到一个小量上. 后者由于 $\bm{M}$ 沿物体缓慢变化，因而才是微小的；而在波长充分大时这些相对论项可以变得重要起来.

\section{铁磁体中的自旋波量子\quad 能谱}

把前节得到的方程运用于波的传播上，在这个波动中磁矩密度除围绕自己的平衡值 $\bm{M}_0$ 进动，还进行微振动. 我们将考察在整个体积内 $\bm{M}_0$ 为恒定的单磁畴样品，而且局限于波长远小于样品线度的情况. 这时，介质就可以认为是无限的.

首先研究的问题只涉及交换相互作用，即基于方程（69.11）. 取 $\bm{M}=\bm{M}_0+\bm{m}$，其中 $\bm{m}$ 是小量，于是可略去 $\bm{m}$ 的二次项而使方程线性化；由于绝对值 $M=M_0$，那么在这一近似中 $\bm{m}\perp\bm{M}_0$. 我们得到
\begin{equation}
\dot{\bm{m}}=\frac{|e|}{mc}\,\alpha_{ik}
\frac{\partial^2\bm{m}}{\partial x_i\partial x_k}\times\bm{M}_0
\end{equation}
（这里及以后取 $g=2$）. 对于与坐标和时间的依赖关系为 $\exp[\mathrm{i}(\bm{k}\cdot\bm{r}-\omega t)]$ 的 $\bm{m}$，我们得到
\begin{equation}
\mathrm{i}\omega\bm{m}=\frac{|e|}{mc}\,\alpha k^2\,\bm{m}\times\bm{M}_0,
\end{equation}
其中 $\alpha=\alpha(\bm{n})=\alpha_{ik}n_in_k$，$\bm{n}$——波矢 $\bm{k}$ 方向上的单位矢量. 这个方程分解成分量，有
\begin{equation*}
\begin{aligned}
\mathrm{i}\omega m_x&=\frac{|e|}{mc}\,\alpha Mk^2m_y,\\
\mathrm{i}\omega m_y&=-\frac{|e|}{mc}\,\alpha Mk^2m_x
\end{aligned}
\end{equation*}

%%p 282
（$z$ 轴沿 $\bm{M}_0$ 方向）. 由此得到自旋波的色散律\footnote{自旋波的平方色散律是布洛赫（F. Bloch, 1930）用微观理论首先得出的. 这个谱的宏观参量的表示式是朗道和栗弗席兹给出的（1945）.}
\begin{equation}
\omega=\frac{|e|M}{mc}\,\alpha(\bm{n})k^2,
\end{equation}
我们看到，如前节初所述，在 $k\to0$ 时，在交换近似下频率趋于零. 自旋波中矢量 $\bm{m}$ 在平面 $xy$ 上以恒定角速度 $\omega$ 旋转，并保证绝对值不变.

在量子力学图像中用公式(70.3)决定自旋波量子的能谱 $\bm{\varepsilon}=\hbar\bm{\omega}$\footnote{在本章中 $\beta$ 处处代表玻尔磁子：$\beta=|e|\hbar/2mc$.}：
\begin{equation}
\varepsilon(\bm{k})=2\beta M\alpha(\bm{n})k^2.
\end{equation}
在二次量子化形式中，描述铁磁体的宏观量将换成自旋波量子的湮没算符和产生算符表述的各算符. 下面将展示应如何处理自旋波量子的(70.4)式.

对应于经典量 $\bm{M}$，我们引入矢量算符 $\hat{\bm{M}}$，它的分量满足确定的对易规则. 令 $\hat{\bm{S}}(\bm{r})\delta V$ 是在 $\bm{r}$ 点物理无限小体元 $\delta V$ 中原子的总自旋算符. 不同体元 $\delta V_1,\delta V_2$ 的算符 $\hat{\bm{S}}(\bm{r}_1)\delta V_1$ 和 $\hat{\bm{S}}(\bm{r}_2)\delta V_2$ 是对易的．同一算符 $\hat{\bm{S}}(\bm{r})\delta V$ 的各分量满足角动量的一般对易规则：
\begin{equation*}
\hat{S}_x\delta V\cdot\hat{S}_y\delta V-\hat{S}_y\delta V\cdot\hat{S}_x\delta V
=\mathrm{i}\hat{S}_z\delta V,
\end{equation*}
或 $\hat{S}_x\hat{S}_y-\hat{S}_y\hat{S}_x=\mathrm{i}\hat{S}_z/\delta V$（其余的对易子也类似）. 在 $\delta V\to0$ 的极限情况下，对任意 $\bm{r}_1$ 和 $\bm{r}_2$，这些规则写成统一的形式
\begin{equation*}
\hat{S}_x(\bm{r}_1)\hat{S}_y(\bm{r}_2)-\hat{S}_y(\bm{r}_2)\hat{S}_x(\bm{r}_1)
=\mathrm{i}\hat{S}_z(\bm{r}_1)\delta(\bm{r}_1-\bm{r}_2).
\end{equation*}
现在，以 $4\beta^2$ 乘这个等式，再注意到磁化强度算符 $\hat{\bm{M}}=-2\beta\hat{\bm{S}}$，便得到
\begin{equation}
\hat{M}_x(\bm{r}_1)\hat{M}_y(\bm{r}_2)-\hat{M}_y(\bm{r}_2)\hat{M}_x(\bm{r}_1)
=-2\mathrm{i}\beta\hat{M}_z(\bm{r}_1)\delta(\bm{r}_1-\bm{r}_2).
\end{equation}
在应用于自旋波时，$\bm{M}$ 描述的是绕 $z$ 轴的微小振动，在小量 $m_x,m_y$ 的一级近似下可以用数量 $M_z\approx M$ 来代换(70.5)式右方的算符 $\hat{M}_z$；这时
\begin{equation}
\hat{m}_x(\bm{r}_1)\hat{m}_y(\bm{r}_2)-\hat{m}_y(\bm{r}_2)\hat{m}_x(\bm{r}_1)
=-2\mathrm{i}\beta M\delta(\bm{r}_1-\bm{r}_2)
\end{equation}
由此可见，量 $m_y$ 和 $m_x$ 在这种情况下（精确到一常数因子）起着正则共轭“广义坐标和广义动量”的作用，这与 $\varphi$ 和 $\bm{\rho}'$ 在液体中声波量子化时所起的作用相类似（§24）. 但是，我们强调指出二者的主要区别. 声子算符的对易规则(24.7)是精确的，同振动是否微小（即同声子态占有数是否微小）无关. (70.6)式的规则是近似的，只在小量 $\bm{m}$ 的一级近似中才是正确的.

从对易规则(70.6)式以及符合线性方程(70.1)的算符 $\hat{m}_x$ 和 $\hat{m}_y$ 间的关系式出发，可以通过自旋波量子的湮没、产生算符来求出 $\hat{m}_x,\hat{m}_y$ 等算符的表达式，这

%%p 283
和§24 中对于声子所做过（参阅§71的习题4）的相仿.

再来研究自旋波量子的能谱，并进而考虑相对论效应对这个谱的影响. 现在已经必须计及 $\bm{M}$ 振动时产生的磁场 $\bm{H}$. 这个场和 $\bm{m}$ 是同级小量；现在用 $\bm{h}$ 表示它.

麦克斯韦方程(69.10)给出
\begin{equation*}
\bm{k}\times\bm{h}=0,\qquad\bm{k}\cdot\bm{h}=-4\pi\bm{k}\cdot\bm{m}.
\end{equation*}
由此可见，场 $\bm{h}$ 的方向沿着波矢量，并等于
\begin{equation}
\bm{h}=-4\pi(\bm{n}\cdot\bm{m})\,\bm{n},
\end{equation}
把(70.7)式代入(69.5)式的被积式的后两项中，我们得到
\begin{equation}
-\bm{m}\cdot\bm{h}-\frac{h^2}{8\pi}=2\pi(\bm{n}\cdot\bm{m})^2
\end{equation}
（这里去掉 $\bm{M}_0\bm{h}$ 项，由于场 $\bm{h}$ 的有势性，对整个体积的积分变为面积分，因而这一项变为零）；在自旋波中，这部分各向异性的能量有时称为静磁能.

假设铁磁体是单轴的，并且属于“易磁化轴”型，因此 $\bm{M}_0$ 的方向沿晶体的对称轴（$z$ 轴）：$\bm{M}_0=\bm{\nu}M$. 鉴于以后的应用，我们还假设存在平行于 $\bm{\nu}$ 方向的外场 $\mathfrak{H}$；此时，样品应该是以 $\bm{\nu}$ 为轴的圆柱体. 于是物体内部的场 $\bm{H}=\mathfrak{H}+\bm{h}$. 线性化的运动方程（这里写的方程已乘过因子 $\hbar$）为
\begin{equation}
-\mathrm{i}\varepsilon\bm{m}=2\beta M\biggl\{\biggl(\alpha k^2+K+\frac{\mathfrak{H}}{M}\biggr)\bm{\nu}\times\bm{m}-\bm{\nu}\times\bm{h}\biggr\}.
\end{equation}
对于单轴晶体 $\alpha=\alpha_1\sin^2\theta+\alpha_2\cos^2\theta$，其中 $\theta$ 是 $\bm{k}$ 和 $\bm{\nu}$ 之间的夹角.

将(70.7)的 $\bm{h}$ 代入到这里，并把方程写成分量式（其中 $x$ 轴最好选在通过 $\bm{\nu}$ 和 $\bm{n}$ 方向（即 $\bm{\nu}\times\bm{n}$）的平面上）. 由获得的 $m_x$ 和 $m_y$ 的两方程的相容性条件得出色散律
\begin{equation}
\varepsilon(\bm{k})=2\beta M\biggl[\biggl(\alpha k^2+K+\frac{\mathfrak{H}}{M}\biggr)
\biggl(\alpha k^2+K+\frac{\mathfrak{H}}{M}+4\pi\sin^2\theta\biggr)\biggr]^{1/2}.
\end{equation}
我们指出，由于存在 $\sin^2\theta=k_x^2/K^2$ 项，将 $\bm{\varepsilon}(\bm{k})$ 按 $\bm{k}$ 分量的幂展开时便不具有简单的幂的特性；这与磁相互作用的远程性有关.

这里对单轴铁磁体（“易磁化轴”型）导出的(70.10)形表示式，对于立方晶体也是正确的. 这是因为，当矢量 $\bm{M}$ 微小偏离其平衡方向时，各向异性能量的改变在两种情况下有同样的形式. 例如，对于 $K'>0$ 的立方晶体，当 $\bm{M}$ 偏离立方体的棱方向的 $\bm{M}_0$ 时，$\delta U_{\text{异}}$ 只与 $\bm{M}$ 和 $\bm{M}_0$ 之间的夹角 $\vartheta$ 有关，并且 $\delta U_{\text{异}}=K'M^2\vartheta^2$. 把此式同单轴晶体相类似的式子 $\delta U_{\text{异}}=KM^2\vartheta^2/2$ 进行比较，我们发现，为过渡到 $K'>0$ 的立方晶体情形，只要在(70.10)式中将 $K$ 换成 $2K'$ 便足够了. 极易相信，用类似的方法，为过渡到 $K'<0$（$\bm{M}_0$ 的方向沿着立方体的空间对角线）的立方晶体需要将 $K$ 换成 $4|K'|/3$. 我们也注意到，在立方晶体中 $\bm{\alpha}(\bm{n})$ 归

%%p 284
结为常数. 而对“易磁化面”型（$K<0$）的单轴铁磁体就是另一种情形了：当 $\bm{M}$ 偏离 $\bm{M}_0$ 时，变更 $\delta U_{\text{异}}$ 既与 $\bm{M}$ 相对于 $\bm{M}_0$ 的极角有关，也与方位角有关；因此这种情形需要特殊研究——参阅习题.

我们记得，(70.10)只属于能谱的开始部分，这部分的准动量 $k\ll1/a$，因此可以进行宏观研究. 从大 $K$ 方面，但是仍然满足 $(\alpha k^2\gg4\pi,K)$ 这个条件来看，表达式(70.10)归结为
\begin{equation}
\varepsilon(\bm{k})=2\beta M\alpha(\bm{n})k^2+2\beta\mathfrak{H}.
\end{equation}
这里的第一项同纯交换作用的表达式(70.4)一致. 外场只给自旋波量子的能量增添 $2\beta\mathfrak{H}$ 项. 于是，在这种近似中自旋波量子在 $\bm{M}_0$ 上磁矩的投影等于 $-2\beta$. 在物体中激发每一个自旋波量子都要使物体的总磁矩减小 $2\beta$.

在相反的情况下，$k\to0$ 时表达式(70.10)趋于一非零值，$(\mathfrak{H}=0$ 时$)$等于
\begin{equation}
\varepsilon(0)=2\beta MK\biggl(1+\frac{4\pi}{K}\sin^2\theta\biggr)^{1/2}.
\end{equation}
于是，考虑磁的各向异性在自旋波的谱中将导致能隙的出现\footnote{相应的频率 $\omega(0)=\varepsilon(0)/\hbar$ 称为\textbf{铁磁共振}频率.}. 这很自然，因为有各向异性时，甚至磁矩整体进行回转（即 $\bm{k}=0$ 时）也与有限的能量相联系． 我们见到，在小 $k$ 时，尽管相对论效应不大，也会给能谱带来相当大的修正.

自旋波量子的概念，作为元激发是关于物体的弱激发态的，因而是与低温有关的. 因此，关于自旋波量子的公式中一切物质常数（也包括磁化强度 $M$）必须选取 $T=0$ 时的值.

回头再来讨论§69中关于耗散微小的假设. 在量子图像中耗散意味着自旋波量子寿命是有限的，这是由于自旋波量子之间以及同其它准粒子的相互作用所造成的.

如果开始讲自旋波量子之间的相互作用，那么首先应该指出在交换近似下自旋波量子的数目是不变的（每个自旋波量子为 $M_z$ 提供同样的贡献 $-2\beta$，而交换作用使 $M_z$ 保持不变）. 因此在此近似下只能发生散射过程． 但是散射概率随温度的降低而减小，这只不过由于散射体数量的减小而造成的，因此在 $T=0$ 时，无论如何，交换性的衰减都要趋于零. 下面可以见到（§72），在交换近似下一个自旋波量子的状态确实是系统的严格的定态.\footnote{还要指出，在交换近似下两个自旋波量子间的散射截面随其能量的减小而趋于零（参阅§73）. 这一情况，在低温时进一步减少了自旋波量子的交换衰减. 在充分低的温度下，对散射过程来说，相对论效应已很重要了.}

在 $T=0$ 时自旋波量子的衰减只是由其裂变过程引起的，这种过程只在计及相对论性相互作用时才有可能，因此它们的概率较小．此外，在小 $k$ 时由于过程终态的统计权重（相体积）小，裂变概率总是降低的.

%%p 285
自旋波量子与声子的相互作用（在这里，交换相互作用哈密顿量中与晶体形变有关的部分起微扰算符的作用）也要引起自旋波量子的衰减. 在 $T=0$ 时，还可能由自旋波量子产生声子的过程，但是，为此自旋波量子的准动量应该充分大——自旋波量子的速度 $\partial\varepsilon/\hbar\partial k$ 应大于声速（参阅254页的注）. 过程概率小也是由于终态统计权重小的缘故.

最后，在铁磁金属中，自旋波量子（由于跟传导电子的交换相互作用）从费米面下，总可能激发电子．这里，在小 $k$ 时，过程概率小也是由于终态的统计权重小的缘故.

\begin{xiti}

试求“易磁化面”型单轴铁磁体中自旋波量子的能谱（$K<0$）.

{\heiti 解：}平衡磁化强度 $\bm{M}_0$ 位于垂直于晶体对称轴（$z$ 轴）的平面上；我们选 $\bm{M}_0$ 方向作为 $x$ 轴，此时磁矩运动的线性方程的形状是
\begin{equation*}
-\mathrm{i}\varepsilon\bm{m}=2\beta\{\alpha k^2\bm{n}_x\times\bm{m}
-|K|m_x\bm{n}_y-\bm{n}_x\times\bm{h}\}.
\end{equation*}
其中 $\bm{n}_x,\bm{n}_y$ 是沿坐标轴的单位矢量，而矢量 $\bm{m}$ 位于垂直于 $\bm{M}_0$ 的 $yz$ 平面上. 将 (70.7)的 $\bm{h}$ 代入，写出分量方程并使所得方程组的行列式等于零，得到自旋波量子的能谱
\begin{equation*}
\varepsilon(\bm{k})=2\beta M\bigl[\alpha k^2(\alpha k^2+|K|)
+4\pi\sin^2\theta(\alpha k^2+|K|\sin^2\varphi)\bigr]^{1/2},
\end{equation*}
其中 $\theta$ 和 $\varphi$ 是 $\bm{k}$ 相对于 $\bm{M}_0$ 方向的极角和方位角（而且方位角由 $xz$ 面上算起）. 在 $\alpha k^2\gg1$ 时，我们又回到平方型谱(70.4)，而在 $k\to0$ 时，自旋波量子的能量趋于：
\begin{equation*}
\varepsilon(0)=4(\pi|K|)^{1/2}\beta M|\sin\theta\sin\varphi|.
\end{equation*}
而当 $\bm{k}$ 位于对称轴与晶体自发磁化强度所构成的 $xz$ 平面上时，上式变为零. 但是，这里趋近于零是近似的：考虑到各向异性能量中的高阶项时，在 $xy$ 面上也会导致各向异性的出现，从而在所有 $\bm{k}$\footnote{我们记得（参阅第八卷§317），各向异性能量按 $\bm{M}$ 幂的展开实际上是按相对论比值 $v/c$ 的展开（因此与 $\bm{M}$ 是否微小，即与是否接近居里点无关）.} 的方向上将会出现有限的能隙.

\end{xiti}

\section{铁磁体中的自旋波量子\quad 热力学量}

铁磁体中激发的自旋波量子对其热力学量有一定的贡献. 当在 $T\ll T_{\mathrm{c}}$ 意义下的低温时，利用前一节所获得的结果可以计算这一贡献. 实际上，在温度为 $T$ 的热平衡时自旋波量子的主要部分有 $\varepsilon\sim T$ 的能量. 对于平方形谱
\begin{equation}
\varepsilon(\bm{k})=2\beta M\alpha(\bm{n})k^2.
\end{equation}
这就是说，温度 $T\ll T_{\mathrm{c}}$ 时，激发的自旋波量子具有准动量 $k\ll(T_{\mathrm{c}}/\beta M\alpha)^{1/2}$. 利

%%p 286
用对 $\alpha$ 估算的(69.7)式，估算出磁化强度为 $M\sim\beta/a^3$（平均一个原胞中的磁矩约为数个 $\beta$），由此得到 $\alpha k\ll1$，即§70结果的适用条件.

铁磁体热力学量的“自旋波量子”（magnon）部分可做为化学势等于零的理想玻色气体的热力学量来计算. 例如，对自旋波量子的热力学势 $\varOmega$（下标用 mag 表示——译注），有
\begin{equation}
\varOmega_{\mathrm{mag}}=T\int\ln(1-\mathrm{e}^{-\varepsilon/T})
\frac{V\mathrm{d}^3k}{(2\pi)^3}
\end{equation}
［参阅第五卷的(54.4)式］. 由此，自旋波量子对内能\footnote{在化学势 $\mu=0$（因此 $\varPhi=N\mu=0$）时，有 $E=\varPhi+TS-PV=TS+\varOmega$；熵 $S=-\partial\varOmega/\partial T$. 当然，无需利用公式(71.2)也可直接写出(71.3)式.}的贡献为
\begin{equation}
E_{\mathrm{mag}}=\varOmega_{\mathrm{mag}}-T\frac{\partial\varOmega_{\mathrm{mag}}}{\partial T}
=\int\frac{\varepsilon}{\mathrm{e}^{\varepsilon/T}-1}\,\frac{V\mathrm{d}^3k}{(2\pi)^3},
\end{equation}
自旋波量子对自发磁化强度的贡献给出后者随温度的变化. 这一贡献，是通过对外磁场的微商而得的，为
\begin{equation*}
M_{\mathrm{mag}}=M(T)-M(0)=-\frac{1}{V}\,\frac{\partial\varOmega_{\mathrm{mag}}}{\partial\mathfrak{H}}
\bigg|_{\mathfrak{H}\to0}
\end{equation*}
［参阅第八卷(31.4)］. 对(71.2)式微商，我们得到
\begin{equation}
M_{\mathrm{mag}}=-\int\frac{\partial\varepsilon}{\partial\mathfrak{H}}
\bigg|_{\mathfrak{H}\to0}\frac{1}{\mathrm{e}^{\varepsilon/T}-1}\,\frac{\mathrm{d}^3k}{(2\pi)^3}.
\end{equation}
微商 $-(\partial\bm{\varepsilon}/\partial\mathfrak{H})$ 是自旋波量子的本征磁矩.

在温度 $T\gg2\pi\beta M$\footnote{对于典型值 $M=2\times10^3$ G，这个条件给出 $T\gg1$ K.} 时，我们计算积分(71.3—71.4)；这时自旋波量子的谱可以利用极限表达式(71.1). 鉴于积分快速收敛，积分可以扩展到整个 $\bm{k}$ 空间（代替一个倒格胞）. 设 $\alpha$ 是常量（对于立方晶体），且作替换 $\mathrm{d}^3k\to4\pi k^2\mathrm{d}k$，经过显然的代入之后，得到
\begin{equation*}
E_{\mathrm{mag}}=\frac{VT^{5/2}}{4\pi^2A^{3/2}}
\int_0^{\infty}\frac{x^{3/2}\,\mathrm{d}x}{\mathrm{e}^x-1}
=V\frac{T^{5/2}\varGamma(5/2)\zeta(5/2)}{4\pi^2A^{3/2}}.
\end{equation*}
为了简化，这里取标记 $A=2\beta M\alpha$（因此 $\varepsilon=Ak^2$）\footnote{关于这一类型的积分计算参阅第五卷§58.}. 于是对于热容量 $C_{\mathrm{mag}}=\partial E_{\mathrm{mag}}/\partial T$，我们得到
\begin{equation}
C_{\mathrm{mag}}=V\frac{5\varGamma(5/2)\zeta(5/2)}{8\pi^2A^{3/2}}\,T^{3/2}
=0.113\biggl(\frac{T}{A}\biggr)^{\!3/2}V.
\end{equation}
我们指出，这个表达式只给出热容量的自旋波量子部分；此外，晶体热容量还包括通常的声子部分.

回到(71.4)式的积分，根据(70.11)式，自旋波量子的磁矩取值 $-2\beta$. 因此

%%p 287
$T\gg2\pi\beta M$ 时得到
\begin{equation}
M_{\mathrm{mag}}=-\frac{\beta T^{3/2}}{2\pi^2A^{3/2}}
\int_0^{\infty}\frac{x^{1/2}\,\mathrm{d}x}{\mathrm{e}^x-1}.
\end{equation}
由此
\begin{equation}
\begin{aligned}
M(T)&=M(0)-\frac{\beta T^{3/2}\varGamma(3/2)\zeta(3/2)}{2\pi^2A^{3/2}}={}\\
&=M(0)-0.117\,\beta\,(T/A)^{3/2}.
\end{aligned}
\end{equation}
（既然声子本身不带磁矩，自旋波量子的贡献当然就包括了磁化强度的所有变化）. 于是，当温度取 $2\pi\beta M\ll T\ll T_{\mathrm{c}}$ 范围时，自发磁化强度的变化便遵守 $T^{3/2}$ 定律（F. Bloch, 1930）.

在自旋波量子的谱中存在能隙(70.10)，在更低的温度范围内将使 $C_{\mathrm{mag}}$ 和 $M_{\mathrm{mag}}$ 对温度有指数函数的依赖关系. 在 $T\ll\beta KM$ 时
\begin{equation}
C_{\mathrm{mag}},M_{\mathrm{mag}}\propto\exp(-2\beta KM/T).
\end{equation}
指数分子中的量在 $\theta=0$ 和 $\theta=\pi$ 时是能隙的最小值（还可参照习题1）.

如果铁磁体的自发磁化强度在基态等于最大值（所谓额定值），对应于物体的所有原子磁矩都平行，那么施加与磁矩同一方向的外磁场时这个值亦不再变化，就是说在这个方向的磁化率 $\chi$ 等于零.

考虑相对论性相互作用时，自发磁化强度（$T=0$ 时）与其“交换”值相比将减小，于是导致出现不为零的磁化率（T. Holstein, H. Primakoff, 1940）. 虽然这个效应很小，但是计算它却有原则性的意义.

在前面计算热力学量的磁性部分时我们舍去“磁性谐振子”的零点能，在这些量的温度关系中它没有贡献. 零点能对应于自旋波量子态的占有数为 1/2：
\begin{equation*}
E(0)_{\mathrm{mag}}=\int\frac{1}{2}\,\varepsilon(\bm{k})\frac{V\mathrm{d}^3k}{(2\pi)^3}.
\end{equation*}
相应地对于“零点”磁化强度有
\begin{equation}
M(0)=-\int\frac{1}{2}\,\frac{\partial\varepsilon}{\partial\mathfrak{H}}\,
\frac{\mathrm{d}^3k}{(2\pi)^3}.
\end{equation}
$k$ 值大时这个积分发散，就是说积分主要决定于短波自旋波量子（$ka\sim1$），对这样的自旋波量子根本不能进行宏观研究. 但是我们将看到，在相对论效应的影响下磁化强度的变化，决定于自旋波量子的长波能谱区域，并可借助于§70中获得的公式来计算.

为简单起见，我们将研究立方晶体，并在此情况下忽略小的各向异性常数，即将自旋波量子谱(70.10)写成形式
\begin{equation}
\varepsilon(\bm{k})=2\beta\bigl[(bk^2+\mathfrak{H})(bk^2+\mathfrak{H}+4\pi M\sin^2\theta)\bigr]^{1/2}
\end{equation}
其中 $b=\alpha M$；在这个表达式中由于考虑了静磁能而产生的 $4\pi M\cdot\sin^2\theta$ 项对应于

%%p 288
相对论效应. 从(71.9)式减去以 $\varepsilon_{\text{交换}}(\bm{k})=2\beta bk^2-2\beta\mathfrak{H}$ 代替 $\bm{\varepsilon}(\bm{k})$ 的同一积分式，则得到待求的相对论效应影响下磁化强度的变化 $\delta M$：
\begin{equation}
\delta M=-\frac{1}{2}\int\frac{\partial}{\partial\mathfrak{H}}
\bigl[\varepsilon(\bm{k})-\varepsilon_{\text{交换}}(\bm{k})\bigr]
\frac{\mathrm{d}^3k}{(2\pi)^3}.
\end{equation}
这个积分在大 $\bm{k}$ 时已变为收敛的了\footnote{为避免误会，我们指出，不可用此法定义基态能量的修正，因为没对 $\mathfrak{H}$ 取微商而利用自旋波量子的长波能谱表达式时，对 $\varepsilon-\varepsilon_{\text{交换}}$ 的积分发散.}.

为了计算上的方便，首先在 $b$ 一定时取 $\delta M$ 对 $M$ 的微商［为此在(71.10)中已引进记号 $b$］. 简单变换后得到
\begin{equation*}
\frac{\partial\delta M}{\partial M}=-\frac{4\pi^2\beta M}{(2\pi)^3}
\int_0^{\pi}\int_0^{\infty}
\frac{\sin^4\theta\cdot2\pi k^2\,\mathrm{d}k\cdot\sin\theta\,\mathrm{d}\theta}
{(bk^2+\mathfrak{H})^{1/2}(bk^2+\mathfrak{H}+4\pi M\sin^2\theta)^{3/2}}.
\end{equation*}
由于对 $\mathrm{d}k$ 的积分收敛，所以能够扩展到 $\infty$.

$\mathfrak{H}=0$ 时积分容易计算；然后对 $M$ 积分，得到
\begin{equation}
\delta M=-\frac{\sqrt{\pi}\,\beta}{8\alpha^{3/2}}.
\end{equation}
这个值很小：$\delta M/M\sim10^{-6}$.

如果外场很强（$\mathfrak{H}\gg4\pi M$），可在被积式的分母中忽略 $4\pi M\cdot\sin^2\theta$ 项. 此后，计算得出结果
\begin{equation}
\delta M=-\frac{2\pi\beta M^{1/2}}{15\alpha^{3/2}\mathfrak{H}^{1/2}}.
\end{equation}
$\mathfrak{H}\to\infty$ 时，$\delta M$ 趋于零. 这是当然要得到的.

在结束时我们指出，如果我们企图把本节在三维情况所使用的方法，用来研究二维铁磁体的磁化强度对温度的依赖关系，那么（在纯交换近似中）代替(71.6)式我们将得到对数发散积分. 这就是说，在所有 $T\neq0$ 的情况下交换相互作用的二维系统实际上不存在自发磁化. 这与§27中指出的二维玻色液体（以及第五卷§137中的二维晶体）的情形相类似. 系统的能量不依赖于磁矩的方向，导致其表达式中只含有矢量 $\bm{M}$ 的导数；结果也导致使磁化破坏的涨落发散（二维情况）. 考虑与 $\bm{M}$ 的方向有关的相对论性相互作用，就可以使涨落得到稳定并使二维铁磁体的存在成为可能.

\begin{xiti}

1. 计算温度 $T\ll\varepsilon(0)$ 时热力学量的自旋波量子部分.

{\heiti 解：}重要的情形是：具有小准动量 $\bm{k}$，在能隙为最小的方向上，即在 $\theta=0$ 和 $\theta=\pi$ 的附近传播的自旋波量子；这两个 $\theta$ 值给出同样的贡献．例如，在小角度 $\theta$

%%p 289
时，对所要求的精度，有
\begin{equation*}
\varepsilon(\bm{k})=2\beta KM+Ak^2+4\pi\beta M\theta^2.
\end{equation*}
其中，对于立方晶体 $A=2\beta M\alpha$，对于“易磁化轴”型的单轴晶体 $A=2\beta M\alpha_2$. 在我们所讨论的温度时自旋波量子的分布可认为是玻尔兹曼分布（即在被积式的分母中忽略 1），并处处以 $\varepsilon(0)$ 替换指数函数前的因子 $\bm{\varepsilon}(\bm{k})$. 把对 $k$ 和 $\theta$ 的积分扩展到 $\infty$，结果得到
\begin{equation*}
E_{\mathrm{mag}}=V\,\frac{KT^{5/2}}{32\pi^{5/2}A^{3/2}}
\exp\bigg(-\frac{2\beta KM}{T}\bigg).
\end{equation*}
\begin{equation*}
M_{\mathrm{mag}}=-\frac{E_{\mathrm{mag}}}{VKM}.
\end{equation*}
在计算热容量时只需对指数函数因子微商；
\begin{equation*}
C_{\mathrm{mag}}=2\beta KMT^{-2}E_{\mathrm{mag}}.
\end{equation*}

2. 在 $\mathfrak{H}\gg4\pi M,\,T\gg\beta\mathfrak{H}$ 条件下，试确定磁化强度与外场的关系.

{\heiti 解：}在所给的条件下，可以忽略相对论项，并将 $\bm{\varepsilon}(\bm{k})$ 写成(70.11)形式. 微商(71.4)式，得
\begin{equation*}
\frac{\partial M}{\partial\mathfrak{H}}=\frac{4\beta^2}{T}
\int\frac{\mathrm{e}^{\varepsilon/T}}{(\mathrm{e}^{\varepsilon/T}-1)^2}\,
\frac{\mathrm{d}^3k}{(2\pi)^3}.
\end{equation*}
在积分中重要的是小 $\bm{k}$，因此
\begin{equation*}
\frac{\partial M}{\partial\mathfrak{H}}\approx4\beta^2T\int\frac{1}{\varepsilon^2(2\pi)^3}\,\mathrm{d}^3k
=\frac{T}{2\pi^2}\int_0^{\infty}\frac{k^2\,\mathrm{d}k}{(\alpha k^2M_0+\mathfrak{H})^2}.
\end{equation*}
（取 $\alpha=\mathrm{const}$；$M_0$ 取 $M$ 在 $\mathfrak{H}=0$ 时的值），最终得到
\begin{equation*}
\frac{\partial M}{\partial\mathfrak{H}}=\frac{T}{8\pi(\alpha M_0)^{3/2}\mathfrak{H}^{1/2}}.
\end{equation*}
于是，在所研究的条件下 $M-M_0\propto\mathfrak{H}^{1/2}$.

3. 试确定于 $T=0$ 时，在弱磁场中，磁化强度对外场的关系.

{\heiti 解：}把(71.10)式的 $\bm{\varepsilon}(\bm{k})$ 代入(71.11)式，再将积分式(71.11)对 $\mathfrak{H}$ 微商，得到
\begin{equation*}
\frac{\partial M}{\partial\mathfrak{H}}=\int
\frac{4\pi^2\beta M_0^2\sin^4\theta}
{\bigl[(\alpha M_0k^2+4\pi M_0\sin^2\theta+\mathfrak{H})(\alpha M_0k^2+\mathfrak{H})^{3/2}\bigr]}\,
\frac{\mathrm{d}^3k}{(2\pi)^3}.
\end{equation*}
在 $\mathfrak{H}\to0$ 时，对 $\mathrm{d}k$ 的积分在小 $k$ 情况下是对数发散的．因此，如只局限于对数的精确性，可以在分母的第一个因式中取 $k=0,\mathfrak{H}=0$，而在第二个因式中取 $\mathfrak{H}=0$，但同时，在 $k^2\sim\mathfrak{H}/\alpha M$ 时从下面并在 $k^2\sim4\pi/\alpha$ 时从上面截断积分. 于是我们获得
\begin{equation*}
\frac{\partial M}{\partial\mathfrak{H}}=\frac{\beta}{32\sqrt{\pi}\,M_0\alpha^{3/2}}
\ln\frac{4\pi M_0}{\mathfrak{H}}.
\end{equation*}

%%p 290
注意，(71.10)式中忽略了 $K$. 在 $\mathfrak{H}\ll KM$ 时在对数中用 $KM_0$ 替代 $\mathfrak{H}$.

4. 在交换近似中试确定在 $r\gg a$ 距离上磁化强度涨落的空间关联函数.

{\heiti 解：}算符 $\hat{m}_x$ 和 $\hat{m}_y$ 满足(70.6)式的对易规则，它们通过自旋波量子湮没和产生算符可表为（在薛定谔绘景中）
\begin{equation*}
\begin{aligned}
\hat{m}_x(\bm{r})&=(\beta M/V)^{1/2}\sum_k
\bigl(\hat{a}_{\bm{k}}\,\mathrm{e}^{\mathrm{i}\bm{k}\cdot\bm{r}}
+\hat{a}_{\bm{k}}^{+}\,\mathrm{e}^{-\mathrm{i}\bm{k}\cdot\bm{r}}\bigr),\\
\hat{m}_y(\bm{r})&=\mathrm{i}(\beta M/V)^{1/2}\sum_k
\bigl(\hat{a}_{\bm{k}}\,\mathrm{e}^{\mathrm{i}\bm{k}\cdot\bm{r}}
-\hat{a}_{\bm{k}}^{+}\,\mathrm{e}^{-\mathrm{i}\bm{k}\cdot\bm{r}}\bigr).
\end{aligned}
\end{equation*}
借助这些算符我们计算关联函数
\begin{equation*}
\varphi_{ik}(\bm{r})=\frac{1}{2}\bigl\langle
\hat{m}_i(\bm{r}_1)\hat{m}_k(\bm{r}_2)
+\hat{m}_k(\bm{r}_2)\hat{m}_i(\bm{r}_1)\bigr\rangle,
\qquad\bm{r}=\bm{r}_1-\bm{r}_2
\end{equation*}
（下标 $i,k$ 取遍 $x,y$ 值）. 考虑到不为零的对角矩阵元，只有乘积
$\langle\hat{a}_{\bm{k}}^{+}\hat{a}_{\bm{k}}\rangle=n_{\bm{k}}$ 和
$\langle\hat{a}_{\bm{k}}\hat{a}_{\bm{k}}^{+}\rangle=n_{\bm{k}}+1$
（其中 $n_{\bm{k}}$ 是自旋波量子态的占有数），我们得到
\begin{equation*}
\varphi_{ik}(\bm{r})=\delta_{ik}\int2\beta M
\biggl(n_{\bm{k}}+\frac{1}{2}\biggr)\mathrm{e}^{\mathrm{i}\bm{k}\cdot\bm{r}}
\frac{\mathrm{d}^3k}{(2\pi)^3}.
\end{equation*}
被积式直接给出关联函数的傅里叶分量. 式中的常数项可以略去：因为在 $\varphi_{ik}(\bm{r})$ 里与它对应的是 $\delta$ 函数项，而整个的研究只是对于 $r>a$ 的距离． 于是
\begin{equation*}
\varphi_{ik}(\bm{k})=2\beta Mn_{\bm{k}}\delta_{ik}
=2\beta M\bigl[\mathrm{e}^{\varepsilon(\bm{k})/T}-1\bigr]^{-1}\delta_{ik},
\end{equation*}
在经典极限情形，$\varepsilon\ll T$，得到
\begin{equation*}
\varphi_{ik}(\bm{k})=\delta_{ik}\,T/2\alpha k^2.
\end{equation*}
在立方铁磁体中 $\alpha=\mathrm{const}$，于是
\begin{equation*}
\varphi_{ik}(\bm{r})=\delta_{ik}\,T/4\pi\alpha r,\qquad
r\gg(\beta M\alpha/T)^{1/2}.
\end{equation*}

\end{xiti}

\section{自旋哈密顿量}

为了得到准动量在整个变域内（不仅仅在长波极限内）自旋波量子的色散律，自然需要利用铁磁体微观结构更细致的概念.

我们所研究的电介质：是由轨道角动量等于零而自旋 $S$ 不为零的原子所组成的．如果我们不涉及与原子电子壳层的激发有关的高激发态，则可以将系统的哈密顿量按基态原子的电子轨道参量求平均（此时，原子核固定在格点上）. 结果我们得到只含原子总自旋算符的系统的自旋哈密顿量.\footnote{这与描写能级精细结构的单原子哈密顿量的构造相类似，可对比第三卷§72.}

如果只考虑与自旋相对取向有关的交换相互作用，那么哈密顿量中原子自旋矢量算符只会以标量结合的形式出现．研究由下列最简单的哈密顿量
\begin{equation}
\hat{H}_{\text{交换}}=-\frac{1}{2}\sum_{m\neq n}J_{nm}\hat{S}_n\hat{S}_m,
\qquad J_{nm}=J(\bm{r}_n-\bm{r}_m)
\end{equation}

%%p 291
所描述的系统，有重要的方法论意义，其中求和遍及所有的原子；“矢量”下标（整数分量）$m$ 和 $n$ 是格点的编号；$\bm{r}_n$ 是格点的径矢. 数 $J_{nm}$ 称为交换积分（对比第三卷§62的习题）\footnote{用自旋哈密顿量描写交换相互作用是狄拉克引进的（P. A. M. Dirac, 1929）. 哈密顿量(72.1)式是凡·弗莱克（J. H. van Vleck, 1931）引进的；它通常称为海森伯哈密顿量，这是因为与此对应的铁磁体模型首先是海森伯研究的.}. 在对 $m$ 和 $n$ 独立求和时，在(72.1)式的求和中每对原子出现两次，而且，当然 $J_{nm}=J_{mn}$.

在(72.1)式中假定晶格的所有磁性原子都是相同的（每个原胞有一个原子）. 作为这个哈密顿量基础的基本假设，认为晶格中的原子相距足够远. 交换积分决定于两个原子波函数的“重叠”，而且随原子间距的增大而很快地（指数地）减小. 因此对于原子间距较大的系统，相互作用可以认为是成对的，因此在(72.1)式中不存在多于两个原子的自旋算符的乘积项. 可以同样精确地认为，两个原子间的交换相互作用每次只由一对电子（每个电子各属一个原子）来实现. 于是相互作用算符将以电子自旋算符的双线性形式构成，而对原子状态平均之后就成为原子自旋的双线性形式了（C. Herring, 1966）\footnote{在此条件下，(72.1)式的求和当然应该只对相邻原子对进行，但是这不能简化公式的书写. 因此，不必明显计及这些条件.}.

如果交换积分 $J_{mn}>0$，哈密顿量(72.1)所描述的系统就是铁磁体. 我们来确定这样系统的基态能量. 这时设想还存在外磁场 $\mathfrak{H}$，给(72.1)式附加算符 $\mathfrak{H}$
\begin{equation}
\hat{V}=-2\beta\mathfrak{H}\sum_m\hat{S}_{mz}
\end{equation}
（$z$ 轴沿外场方向）. 系统总自旋投影算符 $\sum_m\hat{S}_{mz}$ 无论跟 $\hat{H}_{\text{交换}}$ 或跟 $\hat{V}$ 都是可对易的；因此系统的状态可以按这个量的本征值进行分类.

在铁磁情形，与基态对应的是总自旋投影的最大可能值，它等于 $NS$，其中 $N$ 是系统的原子数（自然，这与有无外场无关. 外场只标出选作 $z$ 轴的方向）. 令 $\chi_0$ 是基态归一化的自旋波函数.

如果每个原子自旋投影都取本身的最大值 $S$，总自旋投影才能达到最大值 $NS$. 因此 $\chi_0$ 同时就是每个算符 $\hat{S}_{nz}$ 的本征函数：
\begin{equation}
\hat{S}_{nz}\chi_0=S\chi_0
\end{equation}
我们引入下边所需要的算符 $\hat{S}_{\pm}=\hat{S}_x\pm\mathrm{i}\hat{S}_y$，它们满足对易关系
\begin{equation}
\hat{S}_{+}\hat{S}_{-}-\hat{S}_{-}\hat{S}_{+}=2\hat{S}_z,\qquad
\hat{S}_z\hat{S}_{\pm}-\hat{S}_{\pm}\hat{S}_z=\pm\hat{S}_{\pm}
\end{equation}
（参阅第三卷(26.12)式）. 它们的矩阵元是
\begin{equation}
\langle S_z|S_{+}|S_z-1\rangle=\langle S_z-1|S_{-}|S_z\rangle
=\sqrt{(S+S_z)(S-S_z+1)}
\end{equation}
［参阅第三卷(27.12)式］；算符 $\hat{S}_{+}$ 使分量 $S_z$ 值增加 1，而 $\hat{S}_{-}$ 使之减少 1. 接着我

%%p 292
们写出
\begin{equation*}
\hat{S}_m\hat{S}_n=\hat{S}_{mz}\hat{S}_{nz}
+\frac{1}{2}(\hat{S}_{m+}\hat{S}_{n-}+\hat{S}_{m-}\hat{S}_{n+}),
\end{equation*}
然后写出
\begin{equation}
\hat{H}=-\frac{1}{2}\sum_{m\neq n}J_{mn}
(\hat{S}_{mz}\hat{S}_{nz}+\hat{S}_{m-}\hat{S}_{n+})
-2\beta\mathfrak{H}\sum_m\hat{S}_{mz}.
\end{equation}
其中利用了对称性 $J_{mn}=J_{nm}$ 以及不同原子的算符的可对易性.

由于算符 $\hat{S}_{n+}$ 只有使 $S_z$ 数增加的跃迁矩阵元，因而对于 $S_z$ 数的最大值的态有
\begin{equation}
\hat{S}_{n+}\chi_0=0
\end{equation}
［也可从矩阵元的显式(72.5)中看到］. 因此把哈密顿量算符(72.6)作用于波函数 $\chi_0$ 上，得到
\begin{equation*}
\hat{H}\chi_0=\biggl\{-\frac{1}{2}\sum_{m\neq n}J_{mn}S^2
-2\beta\mathfrak{H}NS\biggr\}\chi_0,
\end{equation*}
花括号中的表达式就是基态能量 $E_0$. 用对 $m$ 和对 $\bm{q}=\bm{n}-\bm{m}$ 的求和代替对 $m$ 和 $n$ 的求和，最终写出 $E_0$ 的形式为
\begin{equation}
E_0=-\frac{1}{2}NS^2\sum_{\bm{q}\neq0}J_{\bm{q}}-2\beta SN\mathfrak{H},
\end{equation}
系统在这个状态的总磁矩是 $2\beta SN$.

依总自旋投影减少的顺序，系统的下一个状态对应于该投影的 $NS-1$ 值；它相当于激发一个磁矩为 $-2\beta$ 的自旋波量子. 总自旋投影的这个值是波函数
\begin{equation}
(2S)^{-1/2}\hat{S}_{n-}\chi_0
\end{equation}
的状态所拥有的. 算符 $\hat{S}_{n-}$ 作用在此状态上，使一个原子自旋投影减少 1\footnote{注意到 $(\hat{S}_{n-}\chi_0)^{*}(\hat{S}_{n-}\chi_0)=\chi_0^{*}\hat{S}_{n-}^{+}\hat{S}_{n-}\chi_0\equiv\langle S|S_{n+}S_{n-}|S\rangle=\langle S|S_{n+}\,|S-1\rangle\langle S-1|S_{n-}\,|S\rangle=2S$，就容易验算函数(72.9)的归一化系数.}. 但是，这个函数并非系统哈密顿量的本征函数；在这一函数中还没有考虑晶格的平移对称性. 哈密顿量的本征函数应该是所有序号 $n$ 的函数(72.9)的线性叠加. 与我们在§55对周期场中电子的布洛赫函数所进行的讨论同样，为了正确考虑平移对称性，这个线性叠加的形式应该是
\begin{equation}
\chi_{\bm{k}}=(2NS)^{-1/2}\sum_n\mathrm{e}^{-\mathrm{i}\bm{k}\cdot\bm{r}_n}
\hat{S}_{n-}\chi_0
\end{equation}
（$N^{-1/2}$ 是归一化因子）. 常矢量 $\bm{k}$ 正是自旋波量子的准动量.

自旋波量子的能量 $\bm{\varepsilon}(\bm{k})$ 是系统的激发态与基态的能量之差 $E_{\bm{k}}-E_0$. 因此
\begin{equation*}
(\hat{H}-E_0)\chi_{\bm{k}}=\varepsilon(\bm{k})\chi_{\bm{k}}
\end{equation*}

%%p 293
将表达式(72.10)代入这一等式的左方，然后以 $\hat{H}\chi_0$ 代替 $E_0\chi_0$，我们得到
\begin{equation}
\varepsilon(\bm{k})\chi_{\bm{k}}=(2NS)^{-1/2}\sum_n
\mathrm{e}^{\mathrm{i}\bm{k}\cdot\bm{r}_n}
(\hat{H}\hat{S}_{n-}-\hat{S}_{n-}\hat{H})\chi_0.
\end{equation}
把 $\hat{H}$ 写成(72.6)式的形式，并利用对易规则(72.4)，不难计算出这一式中的对易子． 再一次考虑 $J_{mn}$ 的对称性，我们算出
\begin{equation}
\hat{H}\hat{S}_{n-}-\hat{S}_{n-}\hat{H}
=\sum_m{}'J_{mn}(\hat{S}_{mz}\hat{S}_{n-}-\hat{S}_{nz}\hat{S}_{m-})
+2\beta\mathfrak{H}\hat{S}_{n-}.
\end{equation}
最后，把上式代入(72.11)式，注意(72.3)式并重新回到对 $\bm{q}=\bm{n}-\bm{m}$ 的求和，我们获得
\begin{equation*}
\varepsilon(\bm{k})\chi_{\bm{k}}=\biggl\{
S\sum_{\bm{q}\neq0}J_{\bm{q}}(1-\mathrm{e}^{\mathrm{i}\bm{k}\cdot\bm{r}_{\bm{q}}})
+2\beta\mathfrak{H}\biggr\}\chi_{\bm{k}}.
\end{equation*}
花括号中的式子就是所求的自旋波量子能量．由于求和号内表达式的虚部是 $\bm{r}_{\bm{q}}$ 的奇函数，所以求和时它变为零，因此最终得到
\begin{equation}
\varepsilon(\bm{k})=S\sum_{\bm{q}\neq0}J_{\bm{q}}
(1-\cos\bm{k}\cdot\bm{r}_{\bm{q}})+2\beta\mathfrak{H}
\end{equation}
（F. Bloch, 1930）.

在哈密顿量(72.1)所描述的系统中，这一公式给出自旋波量子精确的色散律. 自然地，在小 $\bm{k}$ 的极限情况下它变为平方定律：
\begin{equation}
\varepsilon(\bm{k})=\frac{1}{2}S\,k_ik_k\sum_{\bm{q}\neq0}J_{\bm{q}}\,
\chi_{\bm{q}i}\chi_{\bm{q}k}+2\beta\mathfrak{H}
\end{equation}
所研究系统的居里点 $T_{\mathrm{c}}\sim J$，因此当温度 $T\gg J$ 时系统必然成为顺磁体．在这样的温度时，在一级近似下完全可以忽略原子间的相互作用．在这样近似下系统的磁化率将与原子自旋为 $S$ 的理想气体的磁化率一致，并以公式
\begin{equation}
\chi=\frac{N}{V}\,\frac{4\beta^2S(S+1)}{3T}
\end{equation}
表达（参阅第五卷§52）；这是单位体积的磁化率. 这个表达式是函数 $\chi(T)$ 按 $1/T$ 的幂展开的第一项，展开式的其余诸项已经与原子的相互作用有关；我们来确定下一项.

零场中的磁化率由 $\mathfrak{H}\to0$ 时的微商 $\chi=\partial M/\partial\mathfrak{H}$ 定义，而磁化强度是以自由能的微商计算的：$VM=-\partial F/\partial\mathfrak{H}$. 为了求解提出的问题，必须算出精确到 $1/T^2$ 项的自由能 $F$ 的表达式.

根据公式 $F=-T\ln Z$，其中 $Z$ 是配分函数
\begin{equation*}
Z=\sum_n\mathrm{e}^{-E_n/T}\approx\sum_n
\biggl(1-\frac{E_n}{T}+\frac{E_n^2}{2T^2}-\frac{E_n^3}{6T^3}\biggr).
\end{equation*}

%%p 294
求和遍及系统的所有能级\footnote{自由能下一步的计算与第五卷§73中的计算相当，一直进行到展开式的下一项.}. 所研究系统的能谱中能级的总数是有限的，并且等于原子自旋对于晶格取向的一切可能的组合数．每一自旋有 $2S+1$ 个不同的投影；因此所提到的数是 $(2S+1)^N$. 用字母上方的横线表示简单的算术平均值，我们改写 $Z$ 如下
\begin{equation*}
Z=(2S+1)^N\biggl[1-\frac{1}{T}\,\overline{\!E}
+\frac{1}{2T^2}\,\overline{\!E^2}-\frac{1}{6T^3}\,\overline{\!E^3}\biggr].
\end{equation*}
平均值 $\overline{\!E^m}=\mathrm{Tr}\hat{H}^m/(2S+1)^N$. 依照算符迹的已知性质，迹可以由波函数的任意完备系算出；假如这一函数系对应于原子自旋取向的一切可能的组合．于是，求平均便归结为每一自旋对其方向的独立平均，这时 $\overline{\!E}=0$. 现在在 $Z$ 的对数再按 $1/T$ 的幂展开，以相同的精度，我们得到
\begin{equation}
F=-NT\ln(2S+1)-\frac{1}{2T}\,\overline{\!E^2}
+\frac{1}{6T^2}\,\overline{\!E^3}.
\end{equation}
在这一表达式中我们感兴趣的是含 $\mathfrak{H}^2$ 的项，只有这些项对磁化率有贡献. 略去所有其余的项，并注意到自旋分量的奇次幂在求平均时等于零，则得到
\begin{equation*}
F=-\frac{(2\beta\mathfrak{H})^2}{2T}\sum_n\overline{S_{nx}^2}
-\frac{(2\beta\mathfrak{H})^2}{2T^2}\cdot\frac{1}{2}
\sum_{n\neq m}2J_{mn}\overline{(S_nS_{nz})(S_mS_{mz})},
\end{equation*}
平均值
\begin{equation*}
\overline{S_{nz}S_{nx}}=\overline{S_{nz}S_{ny}}=0,\qquad
\overline{S_{nz}^2}=S(S+1)/3.
\end{equation*}
于是
\begin{equation*}
F=-\frac{2}{3T}\beta^2\mathfrak{H}^2NS(S+1)
-\frac{2}{9T^2}\beta^2\mathfrak{H}^2NS^2(S+1)^2
\sum_{\bm{q}\neq0}J_{\bm{q}},
\end{equation*}
因此，最后得到磁化率
\begin{equation}
\chi=\frac{4\beta^2S(S+1)N}{3TV}
\biggl[1+\frac{S(S+1)}{3T}\sum_{\bm{q}\neq0}J_{\bm{q}}\biggr].
\end{equation}
请注意，方括号中修正项的符号与交换积分的符号有关.

\begin{xiti}

1. 以哈密顿量(72.1)描述系统，在温度 $T\gg J$ 时，试计算其热容量的磁性部分.

{\heiti 解：}热容量按 $1/T$ 的幂展开，其第一项可由自由能(72.16)的 $-\overline{\!E^2}/2T$ 项得出. 用同样的方法对哈密顿量(72.1)的平方求平均，得到
\begin{equation*}
\overline{\!E^2}=\frac{1}{4}\,2\sum_{m\neq n}J_{mn}^2\,
\overline{S_{mi}S_{mk}}\ \overline{S_{ni}S_{nk}}
=3\,\frac{S^2(S+1)^2}{9}\,\frac{N}{2}\sum_{\bm{q}\neq0}J_{\bm{q}}^2
\end{equation*}

%%p 295
（因为 $\overline{S_iS_k}=S(S+1)\delta_{ik}/3$）. 结果，我们得到热容量
\begin{equation*}
C_{\text{磁}}=\frac{NS^2(S+1)^2}{6T^2}\sum_{\bm{q}\neq0}J_{\bm{q}}^2
\end{equation*}
这与第五卷公式(73.4)相对应.

2. 忽略自旋间的相互作用，当 $\beta\mathfrak{H}$ 与 $T$ 的比值任意时，计算顺磁体的磁化强度.

{\heiti 解：}配分函数（对磁场中的一个自旋）
\begin{equation*}
Z=\sum_{S_z=-S}^{S}\exp\bigg(-\frac{2\beta\mathfrak{H}}{T}\,S_z\bigg)
=\frac{\sinh\big[2\beta\mathfrak{H}(S+\frac{1}{2})/T\big]}
{\sinh(\beta\mathfrak{H}/T)},
\end{equation*}
计算其自由能，再对 $\mathfrak{H}$ 取微商，我们得到磁化强度
\begin{equation*}
M=\frac{N}{V}\,T\frac{\partial}{\partial\mathfrak{H}}\ln Z
=\frac{2\beta N}{V}\biggl[\biggl(S+\frac{1}{2}\biggr)
\coth\frac{2\beta\mathfrak{H}(S+1/2)}{T}
-\frac{1}{2}\coth\frac{\beta\mathfrak{H}}{T}\biggr]
\end{equation*}
（L. Brillouin, 1927）. 在 $\beta\mathfrak{H}\ll T$ 时，这个表达式变为(72.15)式. 在相反的极限 $\beta\mathfrak{H}\gg T$ 时，磁化强度按规律
\begin{equation*}
M=\frac{2\beta N}{V}\biggl[S-\exp\biggl(-\frac{2\beta\mathfrak{H}}{T}\biggr)\biggr]
\end{equation*}
趋于额面值.

\end{xiti}

\section{自旋波量子的相互作用}

自旋波量子相互作用对铁磁体热力学量磁性部分的贡献，是有重要的方法论意义的问题．我们记得，在§71中的计算是建立在无相互作用自旋波量子的理想气体概念之上的. 现在就交换的自旋哈密顿量(72.1)所描述的系统来研究这一问题.

考虑到贡献只来自于小比值 $T/T_{\mathrm{c}}$ 的最低阶项，我们便可以只限于自旋波量子的成对相互作用. 这就是说，必须研究系统总自旋投影等于 $NS-2$ 的双自旋波量子态.

与此投影相对应的波函数是
\begin{equation}
\begin{aligned}
&\chi_{nn}=[4S(2S-1)]^{-1/2}\hat{S}_{n-}\hat{S}_{n-}\chi_0\\
&\chi_{mn}=(2S)^{-1}\hat{S}_{m-}\hat{S}_{n-}\chi_0,\qquad m\neq n.
\end{aligned}
\end{equation}
由于不同原子的自旋算符是可对易的，所以 $\chi_{mn}=\chi_{nm}$\footnote{如果自旋 $S=1/2$，则以同一个算符 $\hat{S}_{n-}$ 两次作用于基态波函数 $\chi_0$ 上，使其变为零. 于是，在这种情况下所有“对角的”波函数 $\chi_{nn}\equiv0$.}，容易证实函数(73.1)是以条件 $\chi_{mn}^{*}\chi_{mn}=1$ 归一化的，用验证(72.9)式归一化的同样方式展开乘积，就可

%%p 296
做到这一点． 同样可以证实不同函数 $\chi_{mn}$ 是相互正交的.

函数(73.1)本身不是哈密顿量的本征函数. 系统的双自旋波量子定态波函数应该是 $\chi_{mn}$ 的某种线性叠加，我们写出这一波函数为
\begin{equation}
\chi=\sum_{m\neq n}\frac{1}{\sqrt{2}}\,\psi_{mn}\chi_{mn}
+\sum_n\psi_{nn}\chi_{nn}.
\end{equation}
（由于 $\chi_{mn}$ 和 $\chi_{nm}$ 是同一个函数，则应该取 $\psi_{mn}\equiv\psi_{nm}$）. 系数 $\psi_{mn}$ 的集合是某一表象中的波函数，表象的独立变量是晶格中原子的序号. (73.2)式的第一个求和式中引进因子 $1/\sqrt{2}$ 是为了使模平方 $|\chi|^2$ 等于 $\sum|\psi_{mn}|^2$，在后者的求和中不同的 $\psi_{mn}$ 只出现一次.

与建立单自旋波量子定态波函数的方程(72.11)的方法一样，我们可以得到函数(73.2)式应该满足的类似的方程
\begin{equation}
\begin{aligned}
\mathcal{E}\chi=\sum_{m\neq n}\frac{\psi_{mn}}{2^{3/2}S}
\{\hat{H},\hat{S}_{m-}\hat{S}_{n-}\}\chi_0+{}\\
{}+\sum_n\frac{\psi_{nn}}{2[S(2S-1)]^{1/2}}
\{\hat{H},\hat{S}_{n-}\hat{S}_{n-}\}\chi_0.
\end{aligned}
\end{equation}
此处的 $\mathcal{E}=E-E_0$ 是两个相互作用的自旋波量子的能量（括号$\{\dots\}$是对易子）.

将方程(73.3)右方的各对易子展开. 为此我们指出：
\begin{equation*}
\{\hat{H},\hat{S}_{m-}\hat{S}_{n-}\}\equiv
\{\hat{H},\hat{S}_{m-}\}\hat{S}_{n-}
+\hat{S}_{m-}\{\hat{H},\hat{S}_{n-}\},
\end{equation*}
并利用对易子 $\{\hat{H},\hat{S}_{\pm}\}$ 的表达式(72.12)，然后考虑到对易规则(72.4)，把算符 $\hat{S}_z$ 移到最右边的位置上，在这里 $\hat{S}_z$ 作用在函数 $\chi_0$ 上，并将其乘以 $S$. 结果得到
\begin{equation}
\begin{aligned}
\{\hat{H},\hat{S}_{m-}\hat{S}_{n-}\}\chi_0
=S\sum_l\bigl[&J_{ml}(\hat{S}_{m-}-\hat{S}_{l-})\hat{S}_{n-}+{}\\
{}+J_{nl}&(\hat{S}_{n-}-\hat{S}_{l-})\hat{S}_{m-}\bigr]\chi_0
+\delta_{mn}\sum_lJ_{nl}\hat{S}_{n-}\hat{S}_{l-}\chi_0-{}\\
{}-J_{mn}&\hat{S}_{m-}\hat{S}_{n-}\chi_0
+4\beta\mathfrak{H}\hat{S}_{m-}\hat{S}_{n-}\chi_0.
\end{aligned}
\end{equation}
为了简化公式的书写，未写出求和角标的上下限，求和按全部 $l$ 值进行，然而“对角”的 $J_{ll}=0$.

以后的计算步骤是把(73.4)式代入(73.3)式，并使等式两边相同函数 $\chi_{mn}$ 的系数相等，计算虽然十分繁杂，但却是初等的．结果得出 $\psi_{mn}$ 的下列方程组：
\begin{equation}
\begin{aligned}
(2JS-\mathcal{E})\psi_{mn}
=S\sum_l(J_{lm}\psi_{ln}+J_{ln}\psi_{lm})+J_{mn}\psi_{mn}-{}\\
{}-A_S\bigl[J_{mn}(\psi_{mm}+\psi_{nn})
+2\delta_{mn}\sum_lJ_{lm}\psi_{lm}\bigr],
\end{aligned}
\end{equation}
其中
\begin{equation*}
A_S=S\biggl[1-\biggl(\frac{2S-1}{2S}\biggr)^{\!1/2}\biggr].
\end{equation*}

%%p 297
再引进记号 $J$ 代表求和式 $\sum_lJ_{nl}$，而后者显然与下标 $\bm{n}$ 无关\footnote{当所有 $\psi_{nn}$ 都是任意的时候，对于自旋 $S=1/2$ 这些方程也是正确的. 注意到，在 $S=1/2$ 时所有“对角的”量 $\psi_{nn}$ 完全要从 $m\neq n$ 的方程中消失. 在这种情况下，应该简单地认为不存在 $m=n$ 的方程.}.

我们把这个方程从坐标表象（独立变量是原子的坐标 $\bm{r}_n,\bm{r}_m$）变到动量表象，即取
\begin{equation}
\psi_{mn}=\frac{1}{N}\,\mathrm{e}^{\mathrm{i}\bm{K}\cdot(\bm{r}_m+\bm{r}_n)/2}
\sum_{\bm{k}}\psi(\bm{K},\bm{k})\,\mathrm{e}^{\mathrm{i}\bm{k}\cdot(\bm{r}_m-\bm{r}_n)},
\end{equation}
矢量 $\bm{K}$ 代表两个自旋波量子的合准动量，而 $\bm{k}$ 是它们相对运动的准动量；求和是按体积 $Nv$（$N$ 为晶格原子数，$v$ 为晶格原胞的体积）的晶格所允许的 $N$ 个离散值 $\bm{k}$ 进行的. 和 $\psi_{mn}$ 一起，同样也应将下列交换积分表成傅里叶级数形式：
\begin{equation}
J_{mn}=\frac{1}{N}\sum_{\bm{k}}\mathrm{e}^{\mathrm{i}\bm{k}\cdot(\bm{r}_m-\bm{r}_n)}J(\bm{k}),
\qquad J(\bm{k})=\sum_nJ_{0n}\,\mathrm{e}^{-\mathrm{i}\bm{k}\cdot(\bm{r}_0-\bm{r}_n)}
\end{equation}
［因为 $J_{mn}=J_{nm}$，所以 $J(\bm{k})=J(-\bm{k})$］.

略去简单的中间计算，我们直接引出方程(73.5)变换的最终结果
\begin{equation}
\begin{aligned}
\biggl[\varepsilon\biggl(\frac{\bm{K}}{2}+\bm{k}\biggr)
+\varepsilon\biggl(\frac{\bm{K}}{2}-\bm{k}\biggr)-\mathcal{E}\biggr]\psi(\bm{K},\bm{k})+{}\\
{}+\int U(\bm{K},\bm{k},\bm{k}')\psi(\bm{K},\bm{k}')
\frac{V\mathrm{d}^3k'}{(2\pi)^3}=0,
\end{aligned}
\end{equation}
其中
\begin{equation}
\begin{aligned}
NU(\bm{K},\bm{k},\bm{k}')=A_S\biggl[J\biggl(\frac{\bm{K}}{2}+\bm{k}\biggr)
+J\biggl(\frac{\bm{K}}{2}-\bm{k}\biggr)
+J\biggl(\frac{\bm{K}}{2}+\bm{k}'\biggr)+{}\\
{}+J\biggl(\frac{\bm{K}}{2}-\bm{k}'\biggr)\biggr]
-\frac{1}{2}\bigl[J(\bm{k}-\bm{k}')+J(\bm{k}+\bm{k}')\bigr].
\end{aligned}
\end{equation}
而 $\bm{\varepsilon}(\bm{k})$ 是由公式(72.13)确定的单个自旋波量子的能量；用对一个倒格胞的积分来代替对 $\bm{k}'$ 的求和.

这样一来，关于系统的双自旋波量子态的精确求解问题［在哈密顿量(72.1)的范围内］便归结为：求解一个完全类似于动量表象中双粒子系统薛定谔方程［参阅第三卷(130.4)］的方程. 这时，函数 $\bm{\varepsilon}(\bm{k})$ 相当于粒子的动能，而积分方程的核 $U(\bm{K},\bm{k},\bm{k}')$ 相当于相互作用能为 $U$ 的，从动量 $\bm{k}_1,\bm{k}_2$ 态跃迁（散射）到动量为 $\bm{k}_1',\bm{k}_2'$ 态时的矩阵元，其中
\begin{equation*}
\bm{k}_1=\frac{\bm{K}}{2}+\bm{k},\qquad
\bm{k}_2=\frac{\bm{K}}{2}-\bm{k},\qquad
\bm{k}_1'=\frac{\bm{K}}{2}+\bm{k}',\qquad
\bm{k}_2'=\frac{\bm{K}}{2}-\bm{k}',
\end{equation*}
在这一意义上，最好将 $U(\bm{K},\bm{k},\bm{k}')$ 写成形式
\begin{equation}
\begin{aligned}
NU(\bm{k}_1',\bm{k}_2';\bm{k}_1,\bm{k}_2)
=A_S[J(\bm{k}_1)+J(\bm{k}_2)+J(\bm{k}_1')+J(\bm{k}_2')]-{}\\
{}-\frac{1}{2}[J(\bm{k}_1-\bm{k}_1')+J(\bm{k}_1-\bm{k}_2')].
\end{aligned}
\end{equation}

%%p 298
一般情况下，方程(73.8—73.9)很复杂. 我们只在假定 $S\gg1$ 时计算热力学量的修正. 这一情况之所以简单，是由于自旋波量子能量 $\bm{\varepsilon}(\bm{k})$ 正比于 $S$，而它们的相互作用 $U$ 却与 $S$ 无关［在 $S\gg1$ 时，(73.9)中的系数 $A_S\approx1/4$］. 因此 $U$ 可以看成微扰. 这时，来自自旋波量子相互作用而对热力学势 $\varOmega$ 的修正 $\varOmega_{\text{互}}$ 将简单地由 $U$ 的平均值给出．取“对角矩阵元”
\begin{equation}
U(\bm{k}_1,\bm{k}_2;\bm{k}_1,\bm{k}_2)=\frac{1}{2N}
[J(\bm{k}_1)+J(\bm{k}_2)-J(\bm{k}_1-\bm{k}_2)-J(0)],
\end{equation}
我们就可以对给定准动量的自旋波量子态进行平均．然后，用如下的积分来对自旋波量子的平衡分布进行统计平均
\begin{equation}
\varOmega_{\text{互}}=\int n(\bm{k}_1)n(\bm{k}_2)\,
U(\bm{k}_1,\bm{k}_2;\bm{k}_1,\bm{k}_2)
\frac{V^2\mathrm{d}^3k_1\,\mathrm{d}^3k_2}{(2\pi)^6},
\end{equation}
其中 $n(\bm{k})=[\exp(\varepsilon(\bm{k})/T)-1]^{-1}$ 是玻色分布函数.

在低温时，积分定义于小动量 $\bm{k}_1,\bm{k}_2$ 的区域，与此相应，应该把所有的 $\bm{\varepsilon}(\bm{k})$ 和 $J(\bm{k})$ 都按 $\bm{k}$ 的幂展开. 于是，$\bm{\varepsilon}(\bm{k})$ 便由平方形表达式(72.14)给出. 由于 $J(\bm{k})$ 是 $\bm{k}$ 的偶函数，所以它展开的头几项也是平方型：
\begin{equation*}
J(\bm{k})\approx J(0)+a_{ik}k_ik_k,
\end{equation*}
于是
\begin{equation*}
U(\bm{k}_1,\bm{k}_2;\bm{k}_1,\bm{k}_2)
=\frac{1}{N}\,a_{ik}k_{1i}k_{2k}.
\end{equation*}
但是，把这个对于 $\bm{k}_1$ 和 $\bm{k}_2$ 是奇函数的式子代入(73.12)时，由于按 $\bm{k}_1$ 和 $\bm{k}_2$ 的方向求平均，其积分值等于零.

因此，在展开 $J(\bm{k})$ 时必须考虑四次项，结果在积分(73.12)中函数 $U(\bm{k}_1,\bm{k}_2;\bm{k}_1,\bm{k}_2)$ 便是四次幂的形式，并且这个对于 $\bm{k}_1$ 和对于 $\bm{k}_2$ 都是平方的形式带给积分的贡献不等于零．由于积分迅速收敛，所以可将它扩展到整个 $\bm{k}$ 空间. 作变换 $\bm{k}=\bar{\bm{k}}\sqrt{T}$，就可确立 $\varOmega_{\text{互}}$ 与 $T$ 和 $\mathfrak{H}$ 的关系为
\begin{equation}
\varOmega_{\text{互}}=VT^5f(\mathfrak{H}/T).
\end{equation}
并且 $f(0)$ 和 $f'(0)$ 是有限的．由此得到磁化强度的修正项
\begin{equation}
M_{\text{互}}=-\frac{1}{V}\,\frac{\partial\varOmega_{\text{互}}}{\partial\mathfrak{H}}
\bigg|_{\mathfrak{H}=0}=\mathrm{const}\cdot T^4.
\end{equation}
对于热容量的修正也遵从同样的规律进行\footnote{这些结果（任意自旋的普遍情况）是戴森（F. Dyson, 1956）首先得到的. 在叙述方程(73.5)的推导时，我们大体上遵循 R. J. Boyd, J. Callaway（1965）的工作.}.

我们看到，自旋波量子的相互作用只在 $T/T_{\mathrm{c}}$ 的高级近似中才对热力学量有

%%p 299
修正. 我们记得，磁化强度及热容量的磁性部分的主要项都遵守 $T^{3/2}$ 的规律． 在这些项以及与 $\varOmega_{\text{互}}$ 有关的修正项中还有正比于 $T^{5/2}$ 和 $T^{7/2}$ 的项，这些项产生于自旋波量子能量 $\bm{\varepsilon}(\bm{k})$ 按 $k^2$ 幂展开的后续项.

利用得到的方程还可研究两个自旋波量子的束缚态问题. 这些态是以方程(73.8)的离散（在给定 $\bm{K}$ 时）本征值的形式出现的. 这些本征值 $\mathcal{E}(\bm{K})$ 作为变量 $\bm{K}$ 的函数，乃是系统中新的元激发分支. 但是研究表明，只在 $\bm{K}$ 值充分大时这些态才能存在；因此在低温时，这些态无论如何不会影响铁磁体的热力学量\footnote{参阅 M. Wortis, \textit{Phys.\ Rev.\ 132}, 85（1963）. 所谈的是三维晶格问题. 对于二维和一维情况在任何 $\bm{k}$ 时都存在自旋波量子的束缚态.}.

\begin{xiti}

设 $S\gg1$，试求立方晶格的磁化强度和热容量中与自旋波量子相互作用有关的修正项. 在这种晶格中只对相邻（沿立方晶轴）一对原子的交换积分才不为零.

{\heiti 解：}每个原子有六个最邻近的原子. 按定义(73.7)得到
\begin{equation*}
J(\bm{k})=2J_0(\cos k_xa+\cos k_ya+\cos k_za),
\end{equation*}
其中 $J_0$ 是一对邻近原子的交换积分，而 $a$ 是立方晶格的棱边长. 在小 $\bm{k}$ 时
\begin{equation*}
J(\bm{k})\approx J_0\biggl[2-a^2k^2+\frac{a^4}{12}
(k_x^4+k_y^4+k_z^4)\biggr],
\end{equation*}
由此
\begin{equation*}
U(\bm{k}_1,\bm{k}_2;\bm{k}_1,\bm{k}_2)
=-\frac{a^7J_0}{4V}(k_{1x}^2k_{2x}^2+k_{1y}^2k_{2y}^2+k_{1z}^2k_{2z}^2)
\end{equation*}
（$\bm{k}_1$ 和 $\bm{k}_2$ 的奇次项已舍去）. 自旋波量子的能量［根据(72.14)式］为
\begin{equation*}
\varepsilon(\bm{k})=SJ_0a^2k^2+2\beta\mathfrak{H}
\end{equation*}
(72.12)式的积分计算导致以下的结果：
\begin{equation*}
\frac{M_{\text{互}}}{M}=-\frac{3\pi\zeta(3/2)\zeta(5/2)}{2S^2}
\biggl(\frac{T}{4\pi SJ_0}\biggr)^{\!4},
\end{equation*}
\begin{equation*}
C_{\text{互}}=\frac{15\pi\zeta^2(5/2)N}{S}
\biggl(\frac{T}{4\pi SJ_0}\biggr)^{\!4}
\end{equation*}
（$\zeta$ 表示 $\zeta$ 函数）.

\end{xiti}

\section{反铁磁体中的自旋波量子}

反铁磁体的特点是晶格每个原胞里所有电子的磁矩相互抵消（处在无磁场的平衡态）. 严格说来，磁矩密度是按原胞的整个体积分布的. 但在反铁磁体电介质的晶体中可以相当精确地认为磁矩密度实际上是集中在单个原子上，以致

%%p 300
可以用一定的磁矩来描述每一个原子．这些磁矩在所有的原胞里周期性地重复，造成反铁磁体的磁性亚晶格（magnetic sub-lattice）.

各种反铁磁体在结构上很不一样．关于其磁能谱问题我们可以研究一个典型例子，在每个原胞的等价点上（即对晶体的结晶对称性做任何变换时彼此可以互换的点上）有两个磁性原子. 这些亚晶格的原子形成磁矩的平均密度，它们分别用 $\bm{M}_1$ 和 $\bm{M}_2$ 表示，再引入两个矢量
\begin{equation}
\bm{M}=\bm{M}_1+\bm{M}_2,\qquad
\bm{L}=\bm{M}_1-\bm{M}_2.
\end{equation}
在反铁磁体基态 $\bm{M}=0$，$\bm{L}\neq0$，而对于铁磁体 $\bm{M}\neq0$，$\bm{L}=0$. 我们着重指出两者在基态的重要区别. 在交换近似中，铁磁体处于基态时所有磁性原子的自旋投影具有确定（的最大可能）值 $S_z=S$，它对应于磁化强度 $\bm{M}$ 的额面值. 反铁磁体在基态时，显然亚晶格的磁化强度不可能有自己的额面值，因为每个亚晶格的自旋投影之和不是守恒量（甚至在交换近似中也是如此），因此在定态它没有确定值. 甚至各个原子的自旋投影也没有确定值.

可以类似§69中对铁磁体那样来建立矢量 $\bm{L}$ 和 $\bm{M}$ 的宏观“运动方程”的形式，用以描述 $\bm{L},\bm{M}$ 的长波振动. 磁矩密度 $\bm{M}$ 与铁磁情形有同样的方程(69.1)：
\begin{equation}
\frac{\partial\bm{M}}{\partial t}=\gamma\bm{K},
\end{equation}
此时力矩由公式(69.2)来确定.

如果考虑到存在角动量 $\hbar\bm{S}=\bm{M}/\gamma$，一般说来，就表明系统——这里是自旋系统在旋转，便可得出 $\bm{L}$ 的方程. 根据第五卷热力学公式(26.8)，此旋转的角速度可由自由能对此角动量的微商定出：
\begin{equation*}
\bm{\varOmega}=\biggl(\frac{\partial\tilde{F}}{\partial\hbar\bm{S}}\biggr)_{\!T}.
\end{equation*}
此公式显然可以推广到空间有微弱非均匀的情况：
\begin{equation}
\delta\tilde{F}=\int\bm{\varOmega}\,\delta\hbar\cdot\bm{S}\,\mathrm{d}V
=\frac{1}{\gamma}\int\bm{\varOmega}\cdot\delta\bm{M}\,\mathrm{d}V.
\end{equation}
$\bm{L}$ 随时间的变化就是具有角速度 $\bm{\varOmega}$ 的旋转，描述此旋转的方程为
\begin{equation}
\frac{\partial\bm{L}}{\partial t}=\bm{\varOmega}\times\bm{L}.
\end{equation}
现在来定义力矩 $\bm{K}$. 当系统转动无穷小角 $\delta\bm{\phi}$，两个矢量 $\bm{L}$ 和 $\bm{M}$ 都要变化：
\begin{equation*}
\delta\bm{L}=\delta\bm{\phi}\times\bm{L},\qquad
\delta\bm{M}=\delta\bm{\phi}\times\bm{M}.
\end{equation*}
变分自由能 $\tilde{F}$ 并与确定 $\bm{K}$ 的公式(69.2)进行对比，得
\begin{equation}
\bm{K}=\bm{L}\times\bm{H}_L-\frac{1}{\gamma}\,\bm{M}\times\bm{\varOmega},
\end{equation}
根据

%%p 301
\begin{equation}
\delta\tilde{F}=-\int\bm{H}_L\cdot\delta\bm{L}\,\mathrm{d}V,
\end{equation}
并利用对 $\bm{M}$ 变分的公式(74.3)，在式(74.5)中引入有效场 $\bm{H}_L$，它与反铁磁矢量 $\bm{L}$ 相对应.

结果得出 $\bm{M}$ 方程如下
\begin{equation}
\frac{\partial\bm{M}}{\partial t}=\gamma\bm{L}\times\bm{H}_L+\bm{\varOmega}\times\bm{M}.
\end{equation}
这里我们指出方程(74.4)和(74.7)的一些普遍性质. 不难验证此二方程描述的运动是无能耗的. 事实上，能耗等于
\begin{equation*}
T\frac{\partial S}{\partial t}=\biggl(\frac{\partial\tilde{F}}{\partial t}\biggr)_{\!T}
=-\int\bm{H}_L\cdot\frac{\partial\bm{L}}{\partial t}\,\mathrm{d}V
+\frac{1}{\gamma}\int\bm{\varOmega}\cdot\frac{\partial\bm{M}}{\partial t}\,\mathrm{d}V.
\end{equation*}
此方程根据(74.4)和(74.7)而变为零. 其次，由方程(74.4)显然得出 $\dfrac{\partial\bm{L}^2}{\partial t}=0$，这是必然的，因为矢量 $\bm{L}$ 的长度不变. 最后，用 $\bm{M}$ 乘(74.4)式、用 $\bm{L}$ 乘(74.7)式后再相加，得出 $\dfrac{\partial(\bm{L}\cdot\bm{M})}{\partial t}=0$ 即 $\bm{M}$ 和 $\bm{L}$ 相互垂直.

为了确定有效场 $\bm{H}_L$ 和角速度 $\bm{\varOmega}$，需要建立晶体自由能的形式. 在交换近似中，对于所有磁矩（因而矢量 $\bm{L}$ 和 $\bm{M}$）相对于晶格同时转动，自由能应该保持不变. 原胞内两磁性原子位置除具有结晶学等价性的假定之外，由此也能得出对于交换 $\bm{M}_1$ 和 $\bm{M}_2$，即对于 $\bm{L}\to-\bm{L}$，$\bm{M}\to\bm{M}$ 的变换，需要有一个不变性. 由于自由能的不变性，此时也有 $\bm{H}_L\to-\bm{H}_L$，$\bm{\varOmega}\to\bm{\varOmega}$ 的变换.

重要的是: 在所研究的长波极限，磁矩 $\bm{M}$ 是一个小量. 这一点是清楚的，因为假如矢量 $\bm{L}$ 在空间恒定，在交换情况下任何磁矩当然都不会产生. 因而，我们应该计及的项不超过 $\bm{M}$ 的二次项以及 $\bm{L}$ 的一阶微商项. 满足所提条件的表达式有如下形式：
\begin{equation}
F_{\text{交换}}=\int\biggl(\frac{aM^2}{2}
+\frac{1}{2}\alpha_{ik}\frac{\partial\bm{L}}{\partial x_i}
\frac{\partial\bm{L}}{\partial x_k}\biggr)\mathrm{d}V,
\end{equation}
式中系数 $a$ 为正，这对应于平衡时应有 $\bm{M}=0$. 在(74.8)式中归为分部积分的项可以全部略去. $\dfrac{\partial\bm{M}}{\partial x_i}$ 的二次方项可以不必计及，因为这些项显然比 $M^2$ 小.

对积分式(74.8)进行变分（并完成分部积分），得出
\begin{equation}
\bm{H}_L=\frac{1}{2}\alpha_{ik}\frac{\partial^2\bm{L}}{\partial x_i\partial x_k},
\qquad\bm{\varOmega}=\gamma a\bm{M}.
\end{equation}
最后应该指出，磁矩方程(74.7)可改写成连续性方程(69.11)的形式. 现在，磁矩流密度张量的形式如下
\begin{equation*}
\varPi_{il}=\gamma\alpha_{lk}\biggl(\bm{L}\times\frac{\partial\bm{L}}{\partial x_k}\biggr)_{\!i}.
\end{equation*}

%%p 302
$\bm{L}$ 当然没有连续性方程的形式，因为“反铁磁矩”$\int\bm{L}\,\mathrm{d}V$ 不守恒.

现在来研究磁矩的微振动，为此设
\begin{equation*}
\bm{L}=\bm{L}_0+\bm{l}=\bm{\nu}L_0+\bm{l},\qquad\bm{M}=\bm{m},
\end{equation*}
这里 $\bm{l}$ 和 $\bm{m}$ 都是小量，$\bm{\nu}$ 是矢量 $\bm{L}_0$ 平衡方向的单位矢. 运动方程(74.4)和(74.7)线性化之后具有如下形式
\begin{equation}
\frac{\partial\bm{l}}{\partial t}=L_0\bm{\varOmega}\times\bm{\nu},\qquad
\frac{\partial\bm{M}}{\partial t}=\gamma L_0\bm{\nu}\times\bm{H}_L.
\end{equation}
这里考虑了方程(74.7)右侧第二项恒等于零. $\bm{\varOmega}$ 和 $\bm{H}_L$ 是微小变量构成的线性量，因此简单地将 $\bm{L},\bm{M}$ 改成 $\bm{l},\bm{m}$ 是合理的.

对于单色平面自旋波，运动方程(74.10)给出
\begin{equation}
\begin{aligned}
-\mathrm{i}\omega\bm{l}&=\gamma L_0a\,\bm{m}\times\bm{\nu}\\
-\mathrm{i}\omega\bm{m}&=\gamma L_0\alpha(\bm{n})k^2\,\bm{l}\times\bm{\nu},
\end{aligned}
\end{equation}
式中如 §69 仍取 $\bm{\alpha}(\bm{n})=\alpha_{ik}n_in_k$. $\bm{n}$ 为 $\bm{k}$ 方向的单位矢. 这里取第一方程与 $\bm{\nu}$ 的矢积，得出
\begin{equation}
\gamma L_0a\,\bm{m}=\mathrm{i}\omega\,\bm{l}\times\bm{\nu},
\end{equation}
我们见到，在所研究的长波情况，矢量 $\bm{m}$ 却比 $\bm{l}$ 小. 将此表达式代入第二方程，立即得到下述的自旋波色散律：
\begin{equation}
\omega=\gamma L_0k\{a\alpha(\bm{n})\}^{1/2}.
\end{equation}
于是，反铁磁体内的自旋波频率，因而自旋波量子的能量 $\bm{\varepsilon}=\hbar\bm{\omega}$ 在交换近似中正比于 $k$ 而不似铁磁体内正比于 $k^2$\footnote{反铁磁体的这个色散律首先由 L. Hulthen，（1936）得出的. М.\,И.\,Коганов, В.\,М.\,Цукерник，（1958）用宏观研究法导出亚晶格的磁化强度.}. 方程(74.11)建立了 $\bm{l}$ 和 $\bm{m}$ 间的单值联系. 但 $\bm{l}$ 在垂于 $\bm{\nu}$ 的平面上的两个分量仍是任意的. 这表明自旋波有两个独立的极化方向.

从方程(74.12—74.13)看出 $m\sim k(\alpha/a)^{1/2}l\ll l$. 上面已用过相对小的 $\bm{m}$. 为计及磁的各向异性，需要对晶体对称性提出更加具体的假设. 我们认为晶体是单轴对称的，而且 $\bm{L}_0$ 的平衡方向与对称轴重合\footnote{反铁磁体 $\mathrm{FeCO_3}$ 属于这种类型，它是原胞内有两个 Fe 离子的三角晶格（晶类 $D_{3d}$）. 这两个离子的磁矩都沿三重轴的方向.}. 将此方向选为 $z$ 轴.

由于矢量 $\bm{m}$ 较小，只顾及与矢量 $\bm{l}$ 有关的各向异性就足够了. 与铁磁体不同，这里可以略去在振动时所产生的磁场.

在所做的各假设情况下，各向异性能密度 $U_{\text{异}}=Kl^2/2$，且 $K>0$. 这导致有效场 $\bm{H}_L$ 增加一附加项 $-K\bm{l}$，对于平面波有效场
\begin{equation}
\bm{H}_L=-[\bm{\alpha}(\bm{n})k^2+K]\bm{l}.
\end{equation}

%%p 303
由此可见，当计及各向异性时，在(74.12)式中将 $\alpha k^2$ 换成 $\alpha k^2+K$ 即可得到自旋波的色散律. 结果在 $k\to0$ 自旋波量子的能量将不趋于零而趋于有限值\footnote{频率 $\omega(0)=\varepsilon(0)/\hbar$ 称为\textbf{反铁磁共振}频率.}
\begin{equation}
\varepsilon(0)=\hbar\gamma L_0\sqrt{aK}
\end{equation}
（Ch. Kittel, 1951）. 我们注意到，谱的能隙正比于各向异性常数的平方根（而不像(70.12)中的一次幂）. 由于相对论性效应的微弱性表现在各向异性常数的相对微小上，因而这些效应，一般说来，在反铁磁体比铁磁体内更为重要.

可计算自旋波量子对反铁磁体内能的贡献. 根据公式(71.3).（此等式的右侧应乘以 2，因自旋波量子有两个极化方向.）在下列温区
\begin{equation}
\varepsilon(0)\ll T\ll T_N
\end{equation}
（$T_N$ 为尼尔点——反铁磁性消失的温度）可以使用能谱的(74.12)式，在单轴晶体
\begin{equation*}
\omega(\bm{k})=\gamma L_0a^{1/2}\bigl[\alpha_1(k_x^2+k_y^2)+\alpha_2k_z^2\bigr].
\end{equation*}
计算(71.3)式的积分，可得自旋波量子对热容量的贡献，结果如下：
\begin{equation}
C_{\mathrm{mag}}=V\,\frac{4\pi^2T^3}
{15\gamma^3L_0^3a^{3/2}(\alpha_2\alpha_1^2)^{3/2}\hbar^3}.
\end{equation}
当温度 $T\ll\varepsilon(0)$ 时，自旋波量子对热力学量是指数型的小的贡献.

为了确定反铁磁矢量 $\bm{L}$ 与温度的关系，为物体单位体积的能量附加下述形式的一项
\begin{equation}
-\bm{G}\cdot\bm{L}=-GL_z\approx GL\biggl(l-\frac{l^2}{2L}\biggr),
\end{equation}
这里 $\bm{G}$ 是矢量 $\bm{L}$ 的共轭辅助场，沿其平衡方向施加此场.

取自由能对 $G$ 的微商，即可确定 $L$ 的平衡值：$L=-(1/V)(\partial\tilde{F}/\partial G)_T$（比较磁场情况的公式(54.4)（见第五卷）.）结果得到类似于(71.4)的公式
\begin{equation}
L_{\mathrm{mag}}=L(T)-L(0)=-\int\frac{\partial\varepsilon}{\partial G}
\bigg|_{G\to0}\frac{1}{\mathrm{e}^{\varepsilon/T}-1}\,
\frac{2\,\mathrm{d}^3k}{(2p)^3}.
\end{equation}
场 $G$ 的出现导致自旋波量子的色散律中需要替换 $\alpha(\bm{n})k^2\to\alpha(\bm{n})k^2+G/L$. 微商后得出
\begin{equation*}
\frac{\partial\varepsilon}{\partial G}=\frac{\hbar^2\gamma^2La}{2\varepsilon}.
\end{equation*}
代入(74.19)并积分，对于温区(74.16)，最后得到
\begin{equation*}
L(T)-L(0)=-\frac{T^2}{12\gamma L_0^2a^{1/2}(\alpha_2\alpha_1^2)^{1/2}\hbar}.
\end{equation*}
需要指出，在交换近似，如同铁磁情况，积分(74.19)对二维系统是发散的. 这将破坏反铁磁的长程有序性.

\begin{xiti}

%%p 304
试求“易磁化面”型（$K<0$）单轴反铁磁体自旋波量子的能谱.

{\heiti 解：}现在，平衡反铁磁矢量 $\bm{L}_0$ 处于与晶体对称轴（$z$ 轴）相垂的平面. 选 $\bm{L}_0$ 的方向为 $x$ 轴. 各向异性能 $U_{\text{异}}=|K|(\bm{n}_z\bm{l})^2/2$，这里 $\bm{n}_z$，以及后面提到的 $\bm{n}_y$ 分别是 $z,y$ 轴的单位矢. $\bm{H}_L$ 的表达式中出现附加项 $-|K|(\bm{n}_z\cdot\bm{l})\bm{n}_z$. 于是(74.11)中第一方程与方程(74.12)保持不变，而(74.11)中第二方程具有如下形式
\begin{equation*}
-\mathrm{i}\omega\bm{m}=-\gamma L\alpha(\bm{n})k^2(\bm{l}\times\bm{\nu})
-\gamma L|K|(\bm{n}_z\cdot\bm{l})\bm{n}_y.
\end{equation*}
结果，对矢量 $\bm{l}$ 沿 $y$ 轴极化的自旋波量子，其色散律的形式不变，为(74.13)式；而对 $\bm{l}$ 沿 $z$ 轴极化的，则在公式(74.13)中将 $\alpha\bm{k}^2$ 换成 $\alpha\bm{k}^2+|K|$. 此时，各向异性消除了对极化的简并.

\end{xiti}

% 星号节: 原书节号 74* 不占 75 号——
%   thesection 局部显示 74* + 文末 addtocounter -1
\section*{§74*\quad 自旋哈密顿量的反铁磁态}
\phantomsection
\addcontentsline{toc}{section}{\numberline {§74*}自旋哈密顿量的反铁磁态}

在§72我们见到，在铁磁情况可以精确确定哈密顿量(72.1)的基态能量和元激发——自旋波量子的色散律． 对于反铁磁体不可能有这种精确的解. 在节点上自旋方向相反的两套亚晶格的图像，本质上带有经典特性并破坏了在亚晶格之间交换电子的可能. 于是每个亚晶格总自旋投影便没有确定值，因此也不是描述基态和激发态的好量子数.

然而，当亚晶格每个节点的自旋较大时，铁磁亚晶格彼此镶嵌的图像可做为准经典的正确的零级近似，此时量子效应微小，因而可用微扰论来计算\footnote{当 $S\gg1$ 在§73 末已用过微扰论研究自旋波量子的相互作用.}. 对于实际的反铁磁体不能满足 $S\gg1$ 的条件．但是，求解此问题很有意义．下面所述方法对于研究很多磁学理论问题都是有益的．我们指出，在此模型，只与最近邻有相互作用的这种最简单情况，当交换积分 $J_{nm}$ 为负值时，显然要出现反铁磁性. 事实上，在第一级近似，此表式对经典矢量 $\bm{S}_n$ 简单取极小，便决定了基态. 极小对应的状态，是每个节点最近邻的自旋有相反方向，即属于另一个亚晶格.

我们利用(72.6)型的自旋哈密顿量. 首先将此公式中的算符 $\hat{S}_z,\hat{S}_{+},\hat{S}_{-}$ 表示成遵守玻色对易规则的算符. 将自旋 $S$ 写成如下的自旋算符
\begin{equation*}
\hat{S}_{-}=\sqrt{2S}\,\hat{a}^{+}\biggl(1-\frac{\hat{a}^{+}\hat{a}}{2S}\biggr)^{\!1/2},
\qquad
\hat{S}_{+}=\sqrt{2S}\,\hat{a}\biggl(1-\frac{\hat{a}^{+}\hat{a}}{2S}\biggr)^{\!1/2},
\qquad
\hat{S}_z=S-\hat{a}^{+}\hat{a}.
\tag{74*.1}
\end{equation*}
不难证实，如果算符 $\hat{a}^{+},\hat{a}$ 满足普遍的玻色对易规则
\begin{equation*}
\hat{a}\hat{a}^{+}-\hat{a}^{+}\hat{a}=1,
\end{equation*}

%%p 305
则算符 $\hat{S}_z,\hat{S}_{+},\hat{S}_{-}$ 就会满足正确的对易关系(72.4). 例如有
\begin{equation*}
\{\hat{S}_z,\hat{S}_{+}\}=-\sqrt{2S}\biggl(1-\frac{\hat{a}^{+}\hat{a}}{2S}\biggr)^{\!1/2}
\{\hat{a}^{+},\hat{a}\}\hat{a}=\hat{S}_{+}
\end{equation*}
与(72.4)的第二公式相对应（花括号表示对易子）.

类似地可验证
\begin{equation*}
\hat{S}^2\equiv\hat{S}_{-}\hat{S}_{+}+\hat{S}_z^2+\hat{S}_z=S(S+1).
\end{equation*}
最后，从 (74*.1) 式得出：算符 $\hat{S}_{+}$ 作用到最大可能值 $S_z=S$（即，有 $\hat{a}^{+}\hat{a}=0$）的状态，结果为零（理应如此）. 此与 $\hat{S}_{-}$ 作用到 $S_z=-S$ 的状态（$\hat{a}^{+}\hat{a}=2S$）有相同的结果.

这样一来，公式 (74*.1) 给出自旋算符一个精确的表象，它具有一切所需的性质. 此时算符 $\hat{a}^{+}$ 和 $\hat{a}$ 的意义是在所给晶格的节点上“粒子”的产生算符和湮没算符，其投影 $S_z=-1$\footnote{自旋算符的表象 (74*.1) 是 T. Holstein, H. Primakoff 于 1940 年得出的. 他们首先应用变换 (74*.1) 于磁性理论. 讲过的反铁磁近似理论属于 P. W. Anderson, R. Kubo,（1952）.}.

为了方便，从现在开始给不同亚晶格的节点引入不同的记号. 为第一个亚晶格的节点冠以矢量 $\bm{a}$ 下标，第二个亚晶格用 $\bm{b}$ 下标. 将角动量量子化轴取为亚晶格 $a$ 的磁化强度的方向. 对于高 $S$ 的准经典情况，量子涨落较小，于是磁化强度与其额定值接近（下面将计算对此额定值的修正）. 这意味着，算符 $\hat{a}^{+}\hat{a}$，以及算符 $\hat{a}^{+}$ 和 $\hat{a}$ 本身应看成小量，因而在 (74*.1) 的根式中应将其略去. 结果，(74*.1) 式具有如下形式
\begin{equation*}
\hat{S}_{a-}\approx\sqrt{2S}\,\hat{a}_a^{+},\qquad
\hat{S}_{a+}\approx\sqrt{2S}\,\hat{a}_a,\qquad
\hat{S}_{az}=S-\hat{a}_a^{+}\hat{a}_a.
\tag{74*.2}
\end{equation*}
至于第二套亚晶格，其磁化强度沿负 $z$ 轴的方向．准确地说，它的所有算符也有类似于 (74*.2) 的关系式，只是在 $x'=x,\,y'=-y,\,z'=-z$ 的坐标系中（为使带“撇”的坐标形成右手系，$y$ 轴也变号.）. 显然，带撇的算符与通常的算符联系为
$\hat{S}'_{+}=\hat{S}_{-}$，$\hat{S}'_{-}=\hat{S}_{+}$，$\hat{S}'_z=-\hat{S}_z$. 因此，在 $\bm{b}$ 节点处如引入玻色“粒子”（其自旋投影为 $S_z=+1$）的湮没算符 $\hat{b}_b$，则对第二套亚晶格，代替 (74*.2) 有公式：
\begin{equation*}
\hat{S}_{b-}\approx\sqrt{2S}\,\hat{b}_b,\qquad
\hat{S}_{b+}\approx\sqrt{2S}\,\hat{b}_b^{+},\qquad
\hat{S}_{bz}=-S+\hat{b}_b^{+}\hat{b}_b
\tag{74*.3}
\end{equation*}
将 (74*.2), (74*.3) 式代入哈密顿量的表式(72.6)，只保留不高于产生算符和湮没算符的二次项，需要分别对第一、第二套亚晶格的节点求和. 然而，如果给两套亚晶格以同样方式标记节点号码，即
\begin{equation*}
\bm{r}_a-\bm{r}_{a'}=\bm{r}_b-\bm{r}_{b'},\qquad
\text{当 }\bm{a}-\bm{a}'=\bm{b}-\bm{b}'.
\end{equation*}
则公式的书写方式可以简化，（由于亚晶格的等价性，则 $J_{aa'}=J_{bb'}$）. 然后，第二

%%p 306
套亚晶格节点的指标可以用相同的文字 $\bm{a}$ 来表示. 经过不复杂的计算，给出无外场时的 $\hat{H}$
\begin{equation*}
\begin{aligned}
\hat{H}&=S^2\bigl(J_0^{(2)}-J_0^{(1)}\bigr)
+S\bigl(J_0^{(1)}-J_0^{(2)}\bigr)\sum_a
\bigl(\hat{a}_a^{+}\hat{a}_a+\hat{b}_a^{+}\hat{a}_a\bigr)-{}\\
&\qquad{}-S\sum_{aa'}\bigl[J_{aa'}^{(1)}\bigl(\hat{a}_a^{+}\hat{a}_{a'}
+\hat{b}_a^{+}\hat{a}_{a'}\bigr)
+J_{aa'}^{(2)}\bigl(\hat{a}_a^{+}\hat{b}_{a'}^{+}
+\hat{a}_a\hat{b}_{a'}\bigr)\bigr].
\end{aligned}
\tag{74*.4}
\end{equation*}
这里引入记号 $J_{aa'}^{(1)}$ 代表矩阵元 $J_{aa'}$，两节点 $\bm{a},\bm{a}'$ 属于同一个亚晶格；而 $J_{aa'}^{(2)}$ 则属于不同的亚晶格
\begin{equation*}
J_0^{(1,2)}\equiv\sum_aJ_{aa'}^{(1,2)}
\end{equation*}
现在将哈密顿量写在动量表象，为此如(72.7)，设
\begin{equation*}
\begin{aligned}
J_{aa'}^{(1)}=(2/N)\sum_{\bm{k}}\mathrm{e}^{\mathrm{i}\bm{k}\cdot(\bm{r}_a-\bm{r}_{a'})}J_{\bm{k}}^{(1)},
\qquad
J_{\bm{k}}^{(1)}=\sum_{\bm{k}}\mathrm{e}^{-\mathrm{i}\bm{k}\cdot(\bm{r}_a-\bm{r}_{a'})}J_{aa'}^{(1)},\\
\hat{a}_a=\sqrt{2/N}\sum_{\bm{k}}\mathrm{e}^{\mathrm{i}\bm{k}\cdot\bm{r}_a}\hat{a}_{\bm{k}},
\qquad
\hat{a}_{\bm{k}}=\sqrt{2/N}\sum_a\mathrm{e}^{-\mathrm{i}\bm{k}\cdot\bm{r}_a}\hat{a}_a,\\
\hat{a}_a^{+}=\sqrt{2/N}\sum_{\bm{k}}\mathrm{e}^{\mathrm{i}\bm{k}\cdot\bm{r}_a}\hat{a}_{\bm{k}}^{+},
\qquad
\hat{a}_{\bm{k}}^{+}=\sqrt{2/N}\sum_a\mathrm{e}^{-\mathrm{i}\bm{k}\cdot\bm{r}_a}\hat{a}_a^{+},
\end{aligned}
\tag{74*.5}
\end{equation*}
其他的量也与此类似.（计算求和的节点数现在等于 $N/2$，$N$ 为晶格节点总数）. 上面引入的 $J_0^{(1)}$ 正是 $J_{\bm{k}=0}^{(1)}$.

利用这些公式，容易将 (74*.4) 式化成下面形式
\begin{equation*}
\hat{H}=S^2\bigl(J_0^{(2)}-J_0^{(1)}\bigr)
+\sum_{\bm{k}}\bigl[\bigl(A_{\bm{k}}/2\bigr)
\bigl(\hat{a}_{\bm{k}}^{+}\hat{b}_{-\bm{k}}^{+}+\hat{a}_{\bm{k}}\hat{b}_{-\bm{k}}\bigr)
+\bigl(B_{\bm{k}}/2\bigr)
\bigl(\hat{a}_{\bm{k}}^{+}\hat{a}_{\bm{k}}+\hat{b}_{\bm{k}}^{+}\hat{b}_{\bm{k}}\bigr)\bigr],
\tag{74*.6}
\end{equation*}
式中
\begin{equation*}
A_{\bm{k}}=-2SJ_{\bm{k}}^{(2)},\qquad
B_{\bm{k}}=2S\bigl(J_0^{(1)}-J_0^{(2)}-J_{\bm{k}}^{(1)}\bigr).
\tag{74*.7}
\end{equation*}
哈密顿量 (74*.6) 形式上类似于弱非理想玻色气体的哈密顿量(25.7)，只是 $A_{\bm{k}},B_{\bm{k}}$ 的意义与后者的不同，以及存在两种算符 $\hat{a}_{\bm{k}}$ 和 $\hat{b}_{\bm{k}}$. 类似于(25.8)式的变换，可使 (74*.6) 式对角化. 取
\begin{equation*}
\hat{a}_{\bm{k}}=u_{\bm{k}}\hat{c}_{\bm{k}}+v_{\bm{k}}\hat{d}^{+}_{-\bm{k}},
\qquad
\hat{a}_{\bm{k}}^{+}=u_{\bm{k}}\hat{c}_{\bm{k}}^{+}+v_{\bm{k}}\hat{d}_{-\bm{k}},
\tag{74*.8}
\end{equation*}
置换 $\hat{c}_{\bm{k}}$ 和 $\hat{d}_{\bm{k}}$ 算符，得出不同于 (74*.8) 的 $\hat{b}_{\bm{k}}$ 的表达式：
\begin{equation*}
\hat{b}_{\bm{k}}=u_{\bm{k}}\hat{d}_{\bm{k}}+v_{\bm{k}}\hat{c}^{+}_{-\bm{k}},
\qquad
\hat{b}_{\bm{k}}^{+}=u_{\bm{k}}\hat{d}_{\bm{k}}^{+}+v_{\bm{k}}\hat{c}_{-\bm{k}}
\tag{74*.9}
\end{equation*}
新算符 $\hat{c}_{\bm{k}},\hat{c}_{\bm{k}}^{+}$ 和 $\hat{d}_{\bm{k}},\hat{d}_{\bm{k}}^{+}$ 的意义分别是两种独立极化的自旋波量子的湮没算符，产生算符. 如在§25，给 $u_{\bm{k}},v_{\bm{k}}$ 施加条件 $u_{\bm{k}}^2-v_{\bm{k}}^2=1$，这些算符将满足玻色对易规则. 然而，我们以些许不同于§25的方法使哈密顿量对角化.

如果取
\begin{equation*}
u_{\bm{k}}=\cosh\alpha_{\bm{k}},\qquad
v_{\bm{k}}=\sinh\alpha_{\bm{k}},
\tag{74*.10}
\end{equation*}

%%p 307
上述条件可以恒得满足.

将 (74*.8) (74*.9) 代入 (74*.6) 中. 如在公式 (74*.10) 中的参量 $\alpha_{\bm{k}}$ 定义为
\begin{equation*}
\coth(2\alpha_{\bm{k}})=-B_{\bm{k}}/A_{\bm{k}},
\end{equation*}
便消除关于自旋波量子占有数的非对角项.

最后，哈密顿量的形式成为
\begin{equation*}
\hat{H}=E_0+\sum_{\bm{k}}\varepsilon(\bm{k})
\bigl(\hat{c}_{\bm{k}}^{+}\hat{c}_{\bm{k}}+\hat{d}_{\bm{k}}^{+}\hat{d}_{\bm{k}}\bigr),
\tag{74*.11}
\end{equation*}
式中 $\bm{\varepsilon}(\bm{k})$——准动量为 $\hbar\bm{k}$ 的自旋波量子的能量，它等于
\begin{equation*}
\varepsilon(\bm{k})=\frac{1}{2}\sqrt{B_{\bm{k}}^2-A_{\bm{k}}^2}.
\tag{74*.12}
\end{equation*}
根据 (74*.7) 式 $B_{\bm{k}=0}=A_{\bm{k}=0}$. 将这些量按 $\bm{k}$ 分量的幂展开，开始于平方项. 因此，对于小 $\bm{k}$，能量 $\varepsilon(\bm{k})$ 对 $|\bm{k}|$ 是线性式. 这对应于上一节宏观理论的结果. 还要指出，如在 (74*.11) 式所见，在交换近似中对任意的 $\bm{k}$，都会产生对自旋波量子极化的简并.

基态的能量类似(25.13)，由下式给出
\begin{equation*}
E_0=S^2\bigl(J_0^{(2)}-J_0^{(1)}\bigr)
+\sum_{\bm{k}}\biggl[\varepsilon(\bm{k})-\frac{1}{2}B_{\bm{k}}\biggr].
\tag{74*.13}
\end{equation*}
右侧第二项为量子修正项. 显然，它与经典的第一项之比为 $1/S$.

与一个节点的等于 $S$ 的经典磁化强度相比，量子效应使亚晶格的磁化强度减少. 从 (74*.2) 式得出
\begin{equation*}
\langle S_z\rangle=S-\frac{2}{N}\sum_a
\langle\hat{a}_a^{+}\hat{a}_a\rangle
=S-\frac{2}{N}\sum_{\bm{k}}
\langle\hat{a}_{\bm{k}}^{+}\hat{a}_{\bm{k}}\rangle.
\end{equation*}
借助于 (74*.8) 利用自旋波量子算符表述 $\hat{a}_a^{+},\hat{a}_a$. 当
$\langle\hat{c}_{\bm{k}}^{+}\hat{c}_{\bm{k}}\rangle
=\langle\hat{d}_{\bm{k}}^{+}\hat{d}_{\bm{k}}\rangle=0$，
得出 $T=0$ 时类似于(25.18)的公式：
\begin{equation*}
\langle\hat{a}_{\bm{k}}^{+}\hat{a}_{\bm{k}}\rangle
=\frac{B_{\bm{k}}}{4\varepsilon(\bm{k})}-\frac{1}{2}.
\end{equation*}
将此式代入 $\langle S_z\rangle$ 的公式，并从对 $\bm{k}$ 求和过渡到对 $\nu\,\mathrm{d}^3k/(2\pi)^3$ 求积分：
\begin{equation*}
\langle S_z\rangle=S-\int\biggl[\frac{B_{\bm{k}}}{4\varepsilon(\bm{k})}
-\frac{1}{2}\biggr]\frac{\nu\,\mathrm{d}^3k}{(2\pi)^3},
\tag{74*.14}
\end{equation*}
这里 $\nu$ 是亚晶格原胞体积，按对应于亚晶格 $\bm{k}$ 空间的原胞进行 (74*.14) 式中的积分.（记得在 (74*.5) 的求和是按亚晶格的节点进行的）. 对于在最近邻有相互作用的简立方晶格进行数字计算，得出
\begin{equation*}
\langle S_z\rangle\approx S-0.08.
\end{equation*}
应该指出，在此情况甚至外推至 $S=1/2$，对经典值的修正也显得小了．然而，亚晶格的磁化强度可以明显小于晶格的额定值显得令人失望，在这类晶格对远邻有重要的. 与最近邻有不同符号的相互作用，这已与铁磁性相对应.
